\chapter{Formy różniczkowe, orientacja i kohomologia de Rhama}
\label{ch:formy-de-rham}
\index{forma rozniczkowa@Forma różniczkowa}

W Rozdziale~\ref{ch:algebra-tensorowa} zbudowaliśmy algebrę
zewnętrzną jednej przestrzeni wektorowej. Teraz wykonamy tę
konstrukcję w każdym punkcie rozmaitości i pozwolimy, by jej
współczynniki zmieniały się gładko. Otrzymamy formy różniczkowe.
Różniczkowanie form, ich całkowanie i twierdzenie Stokesa
połączą rachunek różniczkowy z homologią przestrzeni z
Rozdziału~\ref{ch:homologia-przestrzeni}. Na końcu porównamy
obie teorie kohomologii. Nadal nie potrzebujemy metryki,
koneksji ani krzywizny; tym pojęciom poświęcimy późniejszy blok.

\section{Czym jest forma różniczkowa?}
\label{sec:formy-definicje}

\begin{definicja}[Formy i ich wartości]
\index{formy i ich wartosci@Formy i ich wartości}
\label{def:forma-rozniczkowa}
Niech $M$ będzie gładką rozmaitością wymiaru $n$.
\emph{Formą różniczkową stopnia $k$} na $M$ nazywamy
gładkie pole naprzemiennych form $k$-liniowych:
\[
 \omega_p\in\Lambda^k T_p^*M,
 \qquad \omega\in\Omega^k(M)
 :=\Gamma(\Lambda^k T^*M).
\]
Przy $k=0$ jest to funkcja gładka,
$\Omega^0(M)=C^\infty(M)$; dla $k>n$ otrzymujemy
$\Omega^k(M)=0$. Przyjmujemy też $\Omega^k(M)=0$ dla $k<0$.
Na rozmaitości z brzegiem gładkość
oznacza lokalne przedłużenie współczynników do
otwartego otoczenia w $\R^n$.
\end{definicja}

Forma $1$-stopnia mierzy wektor styczny i zwraca liczbę.
Forma $2$-stopnia mierzy \emph{zorientowane pole} zadane
parą wektorów; zamiana wektorów zmienia znak.
Jeśli $(U,x)$ jest mapą $x=(x^1,\ldots,x^n)$, to
\begin{equation}\label{eq:forma-lokalna}
 \omega|_U=\sum_{i_1<\cdots<i_k}
 a_{i_1\ldots i_k}\,\dd x^{i_1}\wedge\cdots\wedge\dd x^{i_k},
 \qquad a_I\in C^\infty(U).
\end{equation}
Tutaj $\dd x^i|_p(v)=v(x^i)$ jest kowektorem z
Rozdziału~\ref{ch:rozmaitosci}. Jedyność współczynników
wynika z bazy algebry zewnętrznej
(Twierdzenie~\ref{thm:wedge-baza}). Zmiana mapy zmienia
współczynniki, lecz nie samą formę.

\begin{przyklad}[Dwie formy na płaszczyźnie]
Na $\R^2$ forma $\alpha=x\,\dd y-y\,\dd x$
przypisuje wektorowi $(u,v)$ liczbę $xv-yu$.
Forma $\omega=\dd x\wedge\dd y$ spełnia
\[
 \omega((u_1,v_1),(u_2,v_2))
 =u_1v_2-v_1u_2.
\]
Jest to wyznacznik: wartość bezwzględna daje pole
równoległoboku, znak zaś jego orientację.
\end{przyklad}

\begin{figure}[!htbp]
\centering
\begin{tikzpicture}[>=Stealth,line cap=round,line join=round]
 \fill[manifoldblue!12] (0,0)--(3.2,.45)--(4.25,2.65)--(1.05,2.2)--cycle;
 \draw[manifoldblue!70!black,thick]
  (0,0)--(3.2,.45)--(4.25,2.65)--(1.05,2.2)--cycle;
 \draw[->,manifoldred,very thick] (0,0)--(3.2,.45)
  node[pos=.78,below=5pt] {$u$};
 \draw[->,manifoldgreen!70!black,very thick] (0,0)--(1.05,2.2)
  node[pos=.8,left=5pt] {$v$};
 \draw[->,manifoldgold!90!black,thick]
  (8:1.45) arc[start angle=8,end angle=61,radius=1.45];
 \node[font=\small,align=left,anchor=west,text width=5.3cm]
  at (5.05,1.42)
  {$(\dd x\wedge\dd y)(u,v)=\det(u,v)>0$.
   Po zamianie $u$ i $v$ ten sam obszar
   dostaje przeciwny znak.};
\end{tikzpicture}
\caption{Forma $\dd x\wedge\dd y$ daje pole ze znakiem
w standardowych współrzędnych płaszczyzny. Ogólna $2$-forma
przypisuje liczbę parze wektorów bez wyboru metryki;
nie musi być euklidesowym polem ich równoległoboku.}
\label{fig:forma-zorientowane-pole}
\end{figure}

\section{Iloczyn zewnętrzny, cofnięcie i różniczka}
\label{sec:formy-dzialania}

\begin{stwierdzenie}[Działania punkt po punkcie]
\label{prop:formy-wedge-pullback}
Dla $\alpha\in\Omega^k(M)$ i $\beta\in\Omega^\ell(M)$
iloczyn $\alpha\wedge\beta\in\Omega^{k+\ell}(M)$
otrzymujemy, biorąc iloczyn zewnętrzny w każdym
$T_p^*M$. Zachodzi $\alpha\wedge\beta
=(-1)^{k\ell}\beta\wedge\alpha$.
Jeżeli $F:M\to N$ jest gładkie, to \emph{cofnięcie}
$F^*:\Omega^k(N)\to\Omega^k(M)$ ma wzór
\begin{equation}\label{eq:pullback-formy}
 (F^*\omega)_p(v_1,\ldots,v_k)
 =\omega_{F(p)}(\dd F_pv_1,\ldots,\dd F_pv_k).
\end{equation}
Spełnia $(G\circ F)^*=F^*G^*$ i
$F^*(\alpha\wedge\beta)=F^*\alpha\wedge F^*\beta$.
\end{stwierdzenie}
\begin{proof}
Gładkość iloczynu w mapie wynika z faktu, że jego
współczynniki są sumami iloczynów gładkich $a_I,b_J$;
przemienność z odpowiednim znakiem jest własnością
algebry zewnętrznej z Rozdziału~\ref{ch:algebra-tensorowa}.
Wzór~\eqref{eq:pullback-formy} jest wieloliniowy i
naprzemienny, a w mapach używa pochodnych współrzędnych
$F$, więc wyznacza formę gładką. Reguła łańcucha
$\dd(G\circ F)_p=\dd G_{F(p)}\circ\dd F_p$
daje własność złożenia. Iloczyn zewnętrzny definiuje
się sumą wartości na przestawionych argumentach; podstawienie
do niej tych samych $\dd F_pv_i$ dowodzi ostatniej równości.
\end{proof}

W szczególności $F^*(\dd y^j)=\dd(y^j\circ F)$.
Jeżeli $F:\R^2\to\R^2$, $F(u,v)=(u^2,v+u)$,
to $F^*(\dd x\wedge\dd y)
=\dd(u^2)\wedge\dd(v+u)=2u\,\dd u\wedge\dd v$.
Współczynnik $2u$ jest wyznacznikiem macierzy Jacobiego
i rejestruje także odwrócenie orientacji przy $u<0$.

\begin{twierdzenie}[Różniczka zewnętrzna]
\label{thm:formy-d}
Istnieje dokładnie jedna rodzina $\R$-liniowych
operatorów $\dd:\Omega^k(M)\to\Omega^{k+1}(M)$
o następujących własnościach:
\begin{enumerate}[leftmargin=*]
\item dla funkcji $f$ forma $\dd f$ jest jej różniczką;
\item $\dd(\alpha\wedge\beta)=\dd\alpha\wedge\beta
 +(-1)^k\alpha\wedge\dd\beta$, gdy $\deg\alpha=k$;
\item $\dd^2=0$.
\end{enumerate}
W mapie jej wzór jest szczególnie prosty:
\begin{equation}\label{eq:zewnetrzna-lokalnie}
 \dd\!\left(\sum_I a_I\,\dd x^I\right)
  =\sum_I\sum_j\frac{\partial a_I}{\partial x^j}
     \dd x^j\wedge\dd x^I,
 \qquad \dd x^I=\dd x^{i_1}\wedge\cdots\wedge\dd x^{i_k}.
\end{equation}
Dla każdego gładkiego $F:M\to N$ zachodzi
$F^*\dd\omega=\dd(F^*\omega)$.
\end{twierdzenie}
\begin{proof}
Najpierw sprawdźmy lokalność dowolnego operatora spełniającego
podane warunki. Jeśli $\omega$ znika w otoczeniu $p$, wybierzmy
gładką funkcję $\chi$, równą $1$ blisko $p$, z nośnikiem
w tym otoczeniu. Taką pojedynczą funkcję buduje się w mapie
z $e^{-1/t}$, jak w dowodzie Twierdzenia
\ref{thm:podzial-jednosci-formy}; nie wymaga to różniczki form.
Ponieważ $\chi\omega=0$, reguła Leibniza daje w $p$
$0=\dd(\chi\omega)_p=(\dd\omega)_p$.
Operator zależy więc tylko od formy w otoczeniu punktu.
Lokalne funkcje i formy możemy pomnożyć przez taką funkcję
odcinającą i przedłużyć przez zero; wynik różniczkowania
blisko $p$ nie zależy od przedłużenia.

Każda forma lokalnie jest sumą wyrazów $a_I\dd x^I$.
Możemy już stosować warunki do lokalnych współrzędnych.
Ponieważ $\dd(\dd x^i)=\dd^2x^i=0$,
reguła Leibniza \emph{wymusza} wzór
\eqref{eq:zewnetrzna-lokalnie}; mamy więc jednoznaczność.
Zdefiniujmy operator tym wzorem. Po powtórnym użyciu
otrzymujemy współczynniki
$(\partial_{x^j}\partial_{x^h}a_I)
\dd x^j\wedge\dd x^h\wedge\dd x^I$.
Wyrazy z $(j,h)$ i $(h,j)$ znoszą się, gdyż pochodne
mieszane są równe, a iloczyn zewnętrzny zmienia znak.
Stąd $\dd^2=0$.

Dla dwóch wyrazów $a\dd x^I$ i $b\dd x^J$
zwykła reguła iloczynu $\dd(ab)=a\dd b+b\dd a$
oraz przestawienie $\dd b$ przez $k$ czynników
$\dd x^I$ dają dokładnie znak $(-1)^k$.
Liniowość rozszerza to na sumy.
Sprawdźmy zgodność na przecięciu map. Cofnięcie przez
zmianę współrzędnych zachowuje $\dd f$ dzięki regule
łańcucha i zachowuje $\wedge$ z poprzedniego
stwierdzenia. Zatem dwa lokalne wzory, obliczone po obu
stronach zmiany współrzędnych, są równe na generatorach
$f,\dd x^i$, a przez regułę Leibniza na wszystkich formach.
Wzory sklejają się w globalny operator. Ten sam argument
dla dowolnego $F$ daje $F^*\dd=\dd F^*$.
\end{proof}

\begin{przyklad}[Rachunek na $\R^3$]
Jeżeli $\alpha=P\,\dd x+Q\,\dd y+R\,\dd z$, to
\[
 \dd\alpha=(Q_x-P_y)\,\dd x\wedge\dd y
 +(R_x-P_z)\,\dd x\wedge\dd z
 +(R_y-Q_z)\,\dd y\wedge\dd z.
\]
Dla $\beta=P\,\dd y\wedge\dd z
+Q\,\dd z\wedge\dd x+R\,\dd x\wedge\dd y$
otrzymujemy $\dd\beta=(P_x+Q_y+R_z)
\dd x\wedge\dd y\wedge\dd z$.
Są to wyrażenia występujące w klasycznych wzorach na rotację
i dywergencję. Utożsamianie tu pól wektorowych z formami
korzysta ze standardowej struktury euklidesowej $\R^3$;
sama różniczka zewnętrzna na $M$ nie wymaga metryki.
\end{przyklad}

\begin{definicja}[Kontrakcja i pochodna Liego]
\index{kontrakcja i pochodna liego@Kontrakcja i pochodna Liego}
\label{def:formy-kontrakcja-lie}
Dla pola wektorowego $X$ operator
$\iota_X:\Omega^k(M)\to\Omega^{k-1}(M)$ określamy przez
\[
 (\iota_X\omega)_p(v_1,\ldots,v_{k-1})
 =\omega_p(X_p,v_1,\ldots,v_{k-1}),
 \qquad \iota_Xf=0.
\]
Jeśli $M$ ma brzeg, w definicji przez dwustronny przepływ
zakładamy, że $X$ jest styczne do $\partial M$.
Jeśli $\Phi_t$ jest lokalnym przepływem $X$ z
Rozdziału~\ref{ch:pola}, kładziemy
$\Lie_X\omega=\left.\frac{\dd}{\dd t}\right|_{t=0}\Phi_t^*\omega$.
\end{definicja}

\begin{twierdzenie}[Wzór Cartana]
\label{thm:cartan-formy}
Dla każdej formy $\omega$ mamy
\begin{equation}\label{eq:cartan-formula}
 \boxed{\ \Lie_X\omega=\dd(\iota_X\omega)
       +\iota_X(\dd\omega)\ }.
\end{equation}
Ponadto $\Lie_X\dd=\dd\Lie_X$ oraz
$\iota_X(\alpha\wedge\beta)=
(\iota_X\alpha)\wedge\beta
+(-1)^{\deg\alpha}\alpha\wedge\iota_X\beta$.
\end{twierdzenie}
Dla pola niestycznego do brzegu prawa strona wzoru Cartana
nadal określa gładką formę. Możemy nią zdefiniować
$\Lie_X\omega$ na brzegu: jest to przedłużenie pochodnej
Liego z wnętrza, niezależne od przedłużeń w mapach.
Nie zakładamy wówczas istnienia dwustronnego przepływu w $M$.
\begin{proof}
Ostatnia równość wynika wprost z definicji iloczynu
zewnętrznego: po wstawieniu $X$ w pierwsze miejsce
może on wejść do pierwszej grupy argumentów lub,
po przejściu przez $\deg\alpha$ argumentów, do drugiej.
Ponieważ $\Phi_t^*$ zachowuje iloczyn zewnętrzny,
$\Lie_X$ jest derywacją stopnia $0$ algebry form.
Prawa strona~\eqref{eq:cartan-formula} również nią jest:
rozwijamy oba składniki na $\alpha\wedge\beta$
regułami dla $\dd$ i $\iota_X$; wyrazy mieszane mają
przeciwne znaki i znoszą się.

Wystarczy zatem porównać obie strony na lokalnych
generatorach $f,\dd x^i$. Dla funkcji lewa strona to
$Xf$, a prawa to $\iota_X\dd f=\dd f(X)=Xf$.
Dla $\dd x^i$ lewa strona wynosi
$\dd(\Lie_Xx^i)=\dd(Xx^i)$, gdyż cofnięcie
komutuje z $\dd$; prawa wynosi
$\dd(\iota_X\dd x^i)+\iota_X\dd^2x^i
=\dd(Xx^i)$. Wzór obowiązuje więc dla
każdej sumy wyrazów $f\,\dd x^I$ i globalnie.
Wreszcie, z $\dd^2=0$ dostajemy
$\dd\Lie_X=\dd\iota_X\dd
=\Lie_X\dd$.
\end{proof}

\begin{przyklad}[Pochodna Liego formy]
Na $\R^2$ weźmy $X=\partial_x$ i $\omega=x\,\dd y$.
Przepływ przesuwa punkty poziomo, więc
$\Phi_t^*\omega=(x+t)\dd y$ oraz
$\Lie_X\omega=\dd y$. Ze wzoru Cartana:
$\iota_X\omega=0$, $\dd\omega=\dd x\wedge\dd y$
i $\iota_X\dd\omega=\dd y$. Oba obliczenia są zgodne.
\end{przyklad}
\section{Lokalne istnienie funkcji pierwotnej}
\label{sec:formy-poincare}

Równość $\dd^2=0$ mówi, że każda forma postaci
$\dd\eta$ jest \emph{zamknięta}. Odwrotność nie zachodzi
globalnie: przeszkodę może stanowić topologia. Lokalnie
przeszkoda znika, co wyjaśnia następujący rachunek.
Dla zamkniętej $1$-formy $\omega$ szukamy \emph{funkcji
pierwotnej} $f$ spełniającej $\omega=\dd f$.
Dla formy stopnia $k>1$ jej odpowiednikiem jest
\emph{forma pierwotna} $\eta$ stopnia $k-1$ z
$\omega=\dd\eta$. Lemat Poincarégo obejmuje
oba przypadki.

\begin{lemat}[Homotopia form]
\label{lem:homotopia-form}
Niech $H:[0,1]\times M\to N$ będzie gładką homotopią
odwzorowań $H_0,H_1$. Przy końcach przedziału, a także
przy narożach iloczynu, gdy $M$ ma brzeg, gładkość
rozumiemy jako gładkie przedłużanie składowych w mapach. Istnieje operator liniowy
$K:\Omega^k(N)\to\Omega^{k-1}(M)$ taki, że
\begin{equation}\label{eq:homotopia-form}
 H_1^*\omega-H_0^*\omega
 =\dd(K\omega)+K(\dd\omega).
\end{equation}
\end{lemat}
\begin{proof}
Dla $k=0$ przyjmujemy $K\omega=0$; poniższy rachunek
obejmuje ten przypadek z $a=0$. Każdą formę na
$[0,1]\times M$ zapisujemy jednoznacznie
w postaci $H^*\omega=\dd t\wedge a(t)+b(t)$,
gdzie $a(t)$ jest rodziną form stopnia $k-1$ na $M$,
a $b(t)$ rodziną form stopnia $k$. Położenie
$K\omega:=\int_0^1 a(t)\,\dd t$ oznacza całkowanie
\emph{współczynników} formy po parametrze; można je
wykonać w każdej mapie, a zgodność na przecięciach
wynika z liniowości zmiany współrzędnych.
Obliczenie różniczki daje
\[
 \dd(H^*\omega)
 =\dd t\wedge(\partial_tb-\dd_Ma)+\dd_Mb.
\]
Ponieważ $H^*\dd\omega=\dd H^*\omega$, współczynniki
przy $\dd t$ mówią, że
$a_{\dd\omega}(t)=\partial_tb(t)-\dd_Ma(t)$.
Po scałkowaniu od $0$ do $1$ otrzymujemy
$K\dd\omega=b(1)-b(0)-\dd_MK\omega$.
Ale $b(i)=H_i^*\omega$ dla $i=0,1$,
więc jest to dokładnie~\eqref{eq:homotopia-form}.
\end{proof}

\begin{twierdzenie}[Lemat Poincarégo]
\label{thm:lemat-poincare-formy}
Niech $U\subset\R^n$ będzie otwarty i gwiaździsty
względem $0$. Jeżeli $k\geq1$ i
$\omega\in\Omega^k(U)$ spełnia $\dd\omega=0$,
to istnieje $\eta\in\Omega^{k-1}(U)$ z
$\omega=\dd\eta$. Każda zamknięta forma stopnia
$0$ jest stała na $U$.
\end{twierdzenie}
\begin{proof}
W Lemacie~\ref{lem:homotopia-form} bierzemy
$H(t,x)=tx$. Dla $k\geq1$ odwzorowanie $H_0$
jest stałe, więc $H_0^*\omega=0$; dla zamkniętej
$\omega$ drugi składnik po prawej stronie
\eqref{eq:homotopia-form} także znika.
Dostajemy $\omega=\dd(K\omega)$.
Jeśli chcemy \emph{obliczyć} formę pierwotną
(dla $k=1$: funkcję pierwotną), to
\begin{equation}\label{eq:poincare-pierwotna}
 (K\omega)_x(v_1,\ldots,v_{k-1})
 =\int_0^1 t^{k-1}
   \omega_{tx}(x,v_1,\ldots,v_{k-1})\,\dd t.
\end{equation}
Istotnie, $\dd H(\partial_t)=x$, podczas gdy
$\dd H(v_i)=tv_i$, skąd $k-1$ czynników $t$.
Gwiaździstość zapewnia, że cały odcinek $tx$ leży w $U$.
Dla $k=0$ warunek $\dd f=0$ mówi, że wszystkie
pochodne kierunkowe znikają. Zatem pochodna
$t\mapsto f(tx)$ jest zerowa i $f(x)=f(0)$.
\end{proof}

\begin{przyklad}[Obliczenie funkcji pierwotnej]
Na $\R^2$ forma $\omega=2xy\,\dd x+x^2\,\dd y$
jest zamknięta, bo
$\dd\omega=(2x-2x)\dd x\wedge\dd y=0$.
Wzór~\eqref{eq:poincare-pierwotna} dla $k=1$ daje
\[
 (K\omega)(x,y)=\int_0^1
 \omega_{(tx,ty)}(x,y)\,\dd t
 =\int_0^1 3t^2x^2y\,\dd t=x^2y.
\]
Zatem $f(x,y)=x^2y$ jest funkcją pierwotną
$1$-formy $\omega$, ponieważ $\omega=\dd f$.
Obliczenie pokazuje mechanizm lematu.
\end{przyklad}

\begin{wniosek}[Niezmienniczość przez homotopię]
\label{cor:de-rham-homotopia}
Odwzorowania gładko homotopijne indukują tę samą
operację cofnięcia na kohomologii de Rhama
z Definicji~\ref{def:de-rham-cohomology}.
\end{wniosek}
\begin{proof}
Jeśli $\dd\omega=0$, wzór~\eqref{eq:homotopia-form}
mówi, że $H_1^*\omega-H_0^*\omega=\dd(K\omega)$.
Różnica jest więc dokładna i obie klasy są równe.
\end{proof}

\section{Podział jedności}
\label{sec:formy-podzial-jednosci}

Przy całkowaniu trzeba scalać dane z wielu map.
Służy do tego wynik, który wyjaśnia także wcześniejsze
użycia funkcji odcinających w skrypcie.

\begin{twierdzenie}[Gładki podział jedności]
\label{thm:podzial-jednosci-formy}
Niech $(U_\alpha)$ będzie otwartym pokryciem rozmaitości
$M$ Hausdorffa i o przeliczalnej bazie. Istnieją
nieujemne funkcje gładkie $(\rho_j)$ o lokalnie skończonej
rodzinie nośników (każdy punkt ma otoczenie spotykające
tylko skończenie wiele nośników) oraz indeksy $\alpha(j)$ takie, że
\[
 \supp\rho_j\subset U_{\alpha(j)},
 \qquad \sum_j\rho_j(p)=1\quad(p\in M).
\]
Każdą $\rho_j$ można wybrać o nośniku zwartym
w jednej dziedzinie mapy zawartej w $U_{\alpha(j)}$.
\end{twierdzenie}
\begin{proof}
Najpierw budujemy jedną funkcję odcinającą.
Jeśli w mapie kula domknięta $\overline B_r$
leży w większej kuli $B_R$, przy czym $0<r<R$
i $\overline B_R$ zawiera się w obrazie mapy,
wybierzmy gładką funkcję $\chi$ zależną tylko
od $\|x\|^2$, równą $1$ na $\overline B_r$,
dodatnią między $r$ i $R$, a zerową poza $B_R$.
Można ją złożyć z funkcji $h(t)=e^{-1/t}$ dla
$t>0$ i $h(t)=0$ dla $t\leq0$; jej gładkość w
$t=0$ wynika z zanikania wszystkich pochodnych
$t^{-a}e^{-1/t}$. Przedłużenie $\chi$ przez $0$
na resztę $M$ jest gładkie, bo nośnik ma zwartą
domkniętą otoczkę wewnątrz mapy. Przy brzegu bierzemy
przecięcia tych kul z półprzestrzenią i ograniczamy do nich
te same funkcje. Wszystkie wnętrza i domknięcia poniżej
rozumiemy w topologii $M$, także przy brzegu.

Z lokalnej zwartości i przeliczalności bazy
konstruujemy wyczerpanie zwartymi zbiorami
$K_1\subset\operatorname{int}K_2\subset\cdots$,
$\bigcup_jK_j=M$ (weźmy kolejno skończone
sumy domkniętych, stosunkowo zwartych kul
obejmujące pierwsze elementy przeliczalnego
pokrycia i powiększmy je przy każdym kroku).
Ustalmy $K_j=\varnothing$ dla $j\leq0$.
Zwarty pas $K_j\setminus\operatorname{int}K_{j-1}$
można pokryć skończenie wieloma małymi kulami
w mapach o domknięciach zawartych jednocześnie
w którymś $U_\alpha$ i w
$\operatorname{int}K_{j+1}\setminus K_{j-2}$.
Wybieramy funkcje odcinające dodatnie na
mniejszych kulach pokrywających pas.
Rodzina tych kul dla wszystkich $j$ jest
lokalnie skończona: otoczenie punktu wewnątrz
pewnego $K_N$ nie spotyka kul pochodzących
z pasów o dostatecznie dużym indeksie.

Niech $\chi_j$ oznaczają wszystkie otrzymane
funkcje. Suma $S=\sum_j\chi_j$ jest lokalnie
skończona, gładka i wszędzie dodatnia, bo
każdy punkt należy do któregoś pasa. Kładziemy
$\rho_j=\chi_j/S$. Wtedy nośniki pozostają
w wybranych kulach i
$\sum_j\rho_j=S/S=1$.
Gdy potrzebujemy funkcji indeksowanych samymi zbiorami
pokrycia, sumujemy $\rho_j$ o tym samym $\alpha(j)$.
Nośnik takiej sumy nadal leży w $U_\alpha$:
lokalnie skończona suma zamkniętych nośników jest
zbiorem domkniętym zawartym w $U_\alpha$.
\end{proof}
\section{Orientowalność i orientacje indukowane}
\label{sec:formy-orientacja}

\begin{definicja}[Orientacja]
\index{orientacja@Orientacja}
\label{def:orientacja-formy}
Dwie uporządkowane bazy $m$-wymiarowej przestrzeni
rzeczywistej $V$, dla $m\geq1$, mają tę samą orientację, jeśli
macierz przejścia ma dodatni wyznacznik.
Powstają dwie klasy baz. Na rozmaitości $M^m$
\emph{orientacja} to gładki wybór takiej klasy w
każdym $T_pM$. Równoważnie jest to klasa
nigdzie niezerowych form $m$-stopnia: formy
$\mu,\nu$ wyznaczają tę samą orientację, jeśli
$\nu=f\mu$ dla gładkiej funkcji $f>0$.
Rozmaitość dopuszczająca taki wybór jest
\emph{orientowalna}; wraz z wyborem orientacji
jest \emph{zorientowana}.
\end{definicja}

W wymiarze $0$ stosujemy opis przez formy najwyższego
stopnia: $\Lambda^0 V^*=\R$, a orientacja jest wyborem
jednej z dwóch półprostych, czyli znaku $+1$ lub $-1$.
Każdemu punktowi rozmaitości dyskretnej przypisujemy taki
znak. Opis przez dwie klasy baz dotyczy tylko $m\geq1$,
gdyż przestrzeń zerowa ma tylko jedną bazę pustą.

Dla $m\geq1$ forma $\mu_p$ określa bazę dodatnią $(v_1,\ldots,v_m)$
przez nierówność $\mu_p(v_1,\ldots,v_m)>0$.
To definicja \emph{znaku}, a nie długości ani pola:
zamiana $\mu$ na $2\mu$ nie zmienia orientacji.
Na każdej spójnej orientowalnej rozmaitości istnieją
dwa wybory globalne, wzajemnie przeciwne.

\begin{twierdzenie}[Trzy opisy orientowalności]
\label{thm:orientacja-atlas-forma}
Dla rozmaitości gładkiej $M^m$ równoważne są
poniższe warunki. W mapach przy brzegu dopuszczamy
obie półprzestrzenie $x^m\geq0$ i $x^m\leq0$;
jest to ta sama struktura gładka, a pozwala stosować
odbicie współrzędnej także w wymiarze $1$.
\begin{enumerate}[leftmargin=*]
\item orientowalność $M$;
\item istnienie atlasu, w którym wszystkie wyznaczniki
      odwzorowań przejścia są dodatnie;
\item istnienie globalnej nigdzie niezerowej formy
      $\mu\in\Omega^m(M)$.
\end{enumerate}
\end{twierdzenie}
\begin{proof}
Dla $m=0$ wszystkie trzy warunki są spełnione:
rozmaitość jest dyskretna, dowolny wybór znaków jest
gładki, a funkcja $1$ jest nigdzie niezerową $0$-formą.
Atlas zerowymiarowy sam nie zapisuje wyboru znaków.
Załóżmy więc $m\geq1$.
Forma $\mu$ wybiera dodatnie bazy w każdym punkcie,
więc (3) daje (1). Jeśli wybór (1) jest ustalony,
zmniejszamy dziedziny map do spójnych otoczeń.
Każdą mapę można w razie potrzeby złożyć z odbiciem
jednej współrzędnej, aby jej standardowa baza była
dodatnia. W nowym atlasie dwie bazy lokalne opisują
tę samą orientację, zatem ich wyznacznik przejścia
jest dodatni: (1) daje (2).

Załóżmy (2). W każdej mapie z tego atlasu lokalna
forma $\mu_j=\dd x_j^1\wedge\cdots\wedge\dd x_j^m$
jest dodatnia. Wybierzmy podział jedności $(\rho_j)$
z Twierdzenia~\ref{thm:podzial-jednosci-formy},
o nośnikach w dziedzinach tych map, i przedłużmy
$\rho_j\mu_j$ przez zero. Suma
$\mu=\sum_j\rho_j\mu_j$ jest lokalnie skończona.
W dowolnym punkcie co najmniej jedno $\rho_j$ jest
dodatnie; wszystkie występujące $\mu_j$ są dodatnimi
wielokrotnościami jednej z nich, bo przejścia
mają dodatni wyznacznik. Suma nie może się
wyzerować. Tak dostajemy (3).
\end{proof}

\begin{stwierdzenie}[Dwa podstawowe sposoby przenoszenia orientacji]
\label{prop:orientacja-iloczyn-pullback}
Jeżeli $F:M\to N$ jest dyfeomorfizmem i $N$ ma
orientację $[\nu]$, to $F^*\nu$ orientuje $M$;
$F$ zachowuje orientację wtedy, gdy zgadza się ona
z uprzednio wybraną orientacją $M$. Jeśli $M^m,N^n$
są zorientowane, orientację $M\times N$ wyznacza
$\pr_M^*\mu\wedge\pr_N^*\nu$.
Po zamianie kolejności czynników znak zmienia się
o $(-1)^{mn}$.
\end{stwierdzenie}
\begin{proof}
Różniczka dyfeomorfizmu jest izomorfizmem, więc
cofnięcie nigdzie niezerowej formy najwyższego stopnia
pozostaje nigdzie niezerowe. Na iloczynie wartości
$\mu\wedge\nu$ na bazie najpierw $TM$, potem $TN$
są dodatnie. Przestawienie $m$ wektorów z pierwszego
czynnika z $n$ wektorami drugiego wymaga $mn$
zamian; każda zmienia znak formy.
\end{proof}

Orientacja brzegu i orientacja regularnego zbioru
poziomicowego są dwiema wersjami tej samej zasady:
\emph{najpierw kierunki normalne, potem styczne}.

\begin{definicja}[Orientacja brzegu: normalna na zewnątrz]
\index{orientacja brzegu normalna na zewnatrz@Orientacja brzegu: normalna na zewnątrz}
\label{def:orientacja-brzegu}
Niech $M^m$ będzie zorientowaną rozmaitością z brzegiem.
Dla $m\geq2$ baza $(v_1,\ldots,v_{m-1})$ przestrzeni $T_p\partial M$
jest dodatnia wtedy i tylko wtedy, gdy
$(n,v_1,\ldots,v_{m-1})$ jest dodatnią bazą $T_pM$
dla dowolnego wektora $n$ skierowanego na zewnątrz.
Wymóg nie zależy od konkretnego $n$: dwa takie wektory
mają dodatnio proporcjonalne składowe poprzeczne,
a dodanie wektora stycznego nie zmienia wyznacznika.
Równoważnie, dla formy $\mu$ reprezentującej orientację,
brzeg orientuje $\iota_n\mu$ ograniczona do $T\partial M$.
Ten opis działa również dla $m=1$: znak w punkcie brzegu
jest znakiem liczby $\mu_p(n)$.
\end{definicja}

\begin{stwierdzenie}[Brzeg jest orientowalny]
\label{prop:brzeg-orientowalny}
Powyższa reguła nadaje $\partial M$ gładką orientację.
Dla standardowej orientacji $[a,b]$ końcowy punkt $b$
ma znak $+1$, a początkowy $a$ ma znak $-1$.
Dla dysku $D^2$ ze standardową orientacją płaszczyzny
brzeg biegnie przeciwnie do ruchu wskazówek zegara.
\end{stwierdzenie}
\begin{proof}
W mapie półprzestrzeni $\{x^m\geq0\}$ kierunek
na zewnątrz to $-\partial_{x^m}$; reguła daje
stały wybór znaku baz współrzędnych stycznych.
Zmiany map zachowują orientację $M$ i wysyłają
kierunek na zewnątrz na kierunek na zewnątrz,
więc uzyskane wybory na przecięciach są zgodne.
W wymiarze $1$ przy $b$ wektor zewnętrzny
$+\partial_t$ jest dodatni, a przy $a$ wektor
$-\partial_t$ ujemny. W dysku przy $(1,0)$
wektor zewnętrzny jest $(1,0)$; aby para była
dodatnia, wektor styczny musi wskazywać $(0,1)$.
Ciągłość przenosi ten kierunek na cały okrąg.
\end{proof}

\begin{figure}[!htbp]
\centering
\begin{tikzpicture}[>=Stealth,line cap=round,line join=round]
 \fill[manifoldblue!9] (0,0) circle (1.65);
 \draw[manifoldblue!75!black,line width=1.2pt] (0,0) circle (1.65);
 \draw[->,manifoldred,very thick] (1.65,0)--(2.73,0)
  node[right,font=\small] {$n$};
 \draw[->,manifoldgreen!70!black,very thick] (1.65,0)--(1.65,1.15)
  node[above,font=\small] {$v$ styczny};
 \draw[->,manifoldgold!85!black,thick]
  (210:1.65) arc[start angle=210,end angle=255,radius=1.65];
 \draw[->,manifoldgold!85!black,thick]
  (120:1.65) arc[start angle=120,end angle=165,radius=1.65];
 \fill[manifoldred] (1.65,0) circle(2.5pt)
  node[below=5pt] {$p$};
 \node[font=\small,align=left,anchor=west,text width=4.6cm]
  at (4.75,.55)
  {Para $(n,v)$ ma dodatni wyznacznik.
   Dodatni kierunek na $\partial D^2$
   biegnie przeciwnie do wskazówek zegara.};
\end{tikzpicture}
\caption{Orientację brzegu wyznacza reguła
„normalna na zewnątrz jako pierwszy wektor”.}
\label{fig:orientacja-brzegu-dysk}
\end{figure}

\begin{twierdzenie}[Orientacja regularnego przeciwobrazu]
\label{thm:orientacja-poziomicy}
Niech $F:M^m\to N^q$ będzie gładkie,
$M,N$ zorientowane, $N$ bez brzegu, a $y\in N$ wartością
regularną. Jeśli $M$ ma brzeg, zakładamy ponadto, że $y$
jest wartością regularną $F|_{\partial M}$.
Wtedy $S=F^{-1}(y)$ jest orientowalną podrozmaitością
wymiaru $m-q$ (o ile nie jest pusta). Orientację
można określić dla wszystkich wymiarów następująco.
Dla dodatnich form objętości $\mu_p$ na $T_pM$
i $\nu_y$ na $T_yN$, dowolnej bazy $(e_1,\ldots,e_q)$
oraz podniesień $\dd F_pw_i=e_i$ kładziemy
\[
 \eta_p(v_1,\ldots,v_{m-q})
 =\frac{\mu_p(w_1,\ldots,w_q,v_1,\ldots,v_{m-q})}
        {\nu_y(e_1,\ldots,e_q)}.
\]
Puste listy argumentów oznaczają wartości form stopnia $0$.
Forma $\eta$ określa orientację $S$. Gdy $q\geq1$
i $m-q\geq1$, jest to reguła: $(v_1,\ldots,v_{m-q})$
w $T_pS=\ker\dd F_p$ jest dodatnia, jeśli dla
wektorów $w_1,\ldots,w_q\in T_pM$, których obrazy
$(\dd F_pw_1,\ldots,\dd F_pw_q)$ tworzą dodatnią
bazę $T_yN$, ciąg
$(w_1,\ldots,w_q,v_1,\ldots,v_{m-q})$
jest dodatni w $T_pM$. Gdy $m=q$, punktowi $p\in S$
przypisujemy znak izomorfizmu $\dd F_p$ między
zorientowanymi przestrzeniami. Przy brzegu
$\partial S=S\cap\partial M$.
\end{twierdzenie}
\begin{proof}
Regularność znaczy, że $\dd F_p$ jest surjekcją,
więc takie $w_i$ istnieją. Przy $\partial M$ dodatkowa
surjekcja na przestrzeni stycznej do brzegu pozwala wybrać
$q$ współrzędnych stycznych, w których pochodne $F$ są
niezależne. Twierdzenie o funkcji odwrotnej zastosowane
do $F$ wraz z pozostałymi współrzędnymi, w tym współrzędną
normalną $s\geq0$, daje lokalną postać poziomicy
$\R^{m-q-1}\times[0,\infty)$. Stąd podana formuła brzegu.
Bez założenia o $F|_{\partial M}$ wniosek mógłby zawieść:
dla $F(x,s)=x^2+s$ na $s\geq0$ wartość $0$ jest regularna
dla $F$, lecz przeciwobraz jest pojedynczym punktem,
a nie oczekiwaną jednowymiarową podrozmaitością z brzegiem. Zastąpienie $w_i$
innymi podniesieniami zmienia je o elementy
$\ker\dd F_p$; macierz zmiany bazy ma wówczas
blok trójkątny z jednostką na przekątnej w części
normalnej, zatem nie zmienia znaku. Zmiana
dodatniej bazy w $T_yN$ ma dodatni wyznacznik
i także zachowuje znak. Reguła jest więc
jednoznaczna. W ilorazowym wzorze na $\eta_p$ zmiana
dowolnej bazy $(e_i)$ mnoży licznik i mianownik przez
ten sam niezerowy wyznacznik; zmiana dodatnich form
$\mu_p,\nu_y$ mnoży wynik przez liczbę dodatnią.
W lokalnej postaci submersji z
Wniosku~\ref{wn:submersja-lokalnie} wybory normalnych
i stycznych są gładkie; znaki są zgodne na przecięciach.
To dowodzi orientowalności $S$ i wzoru.
\end{proof}

Ta konstrukcja działa też dla zorientowanej
\emph{wiązki normalnej} podrozmaitości $S\subset M$:
jej baza dodatnia, a następnie baza $TS$ mają razem
dawać bazę dodatnią $TM$. Wybór orientacji normalnej
nazywamy \emph{koorientacją}. Bez takiego wyboru sama
orientacja $M$ nie zawsze orientuje $S$.

\begin{przyklad}[Okrąg jako poziomica]
W $\R^2$ i $\R$ wybierzmy orientacje standardowe
i weźmy $F(x,y)=x^2+y^2$. Na $S^1=F^{-1}(1)$
różniczka ma wartość dodatnią na wektorze
skierowanym promieniowo na zewnątrz.
Twierdzenie~\ref{thm:orientacja-poziomicy} daje więc
orientację przeciwną do ruchu wskazówek zegara,
taką samą jak orientacja brzegu $D^2$.
\end{przyklad}

\begin{przyklad}[Co zawodzi na wstędze i przestrzeni rzutowej]
Wstęga Möbiusa z Przykładu~\ref{prz:mobius}
po jednym okrążeniu odwraca
kierunek poprzeczny, zachowując kierunek
wzdłuż okręgu. Orientacja płaszczyzny stycznej
musiałaby więc zmienić znak i nie może być globalna.
Jej brzeg jest jednak okręgiem, a zatem orientowalnym.

Dla $n\geq1$ przestrzeń $\mathbb{RP}^n$ jest
orientowalna dokładnie wtedy, gdy $n$ jest nieparzyste.
Istotnie, nakrycie $S^n\to\mathbb{RP}^n$ identyfikuje
$x$ z $-x$. Antypodalne $A(x)=-x$ zmienia orientację
$S^n$ o znak $(-1)^{n+1}$: dodatnią bazę przestrzeni otaczającej
otrzymujemy z $(x,v_1,\ldots,v_n)$ w $\R^{n+1}$,
a $A$ mnoży wszystkie $n+1$ wektorów przez $-1$.
Orientacja na ilorazie istnieje wtedy i tylko wtedy,
gdy to utożsamienie zachowuje orientację sfery:
jeśli zachowuje, orientacja schodzi przez lokalne
dyfeomorfizmy; jeśli odwraca, jej podniesienie z
ilorazu byłoby zarazem zachowane i odwrócone, co
jest niemożliwe na spójnej sferze.
\end{przyklad}
\section{Całkowanie form na rozmaitości}
\label{sec:formy-calkowanie}

Na zorientowanej rozmaitości $M^m$ całkujemy formy
stopnia $m$: mają dokładnie tyle argumentów, ile
niezależnych kierunków ma $M$. Bez orientacji można
całkować \emph{gęstości}, ale wymagają osobnej definicji;
całka z formy najwyższego stopnia potrzebuje znaku.

\begin{definicja}[Całka z formy o zwartym nośniku]
\index{calka z formy o zwartym nosniku@Całka z formy o zwartym nośniku}
\label{def:calka-z-formy}
Niech $M^m$ będzie zorientowana i
$\omega\in\Omega_c^m(M)$, gdzie dolny indeks $c$
oznacza zwarty nośnik. W mapie $(U,x)$ zgodnej z
orientacją, jeśli $\supp\omega\subset U$ i
$\omega=f\,\dd x^1\wedge\cdots\wedge\dd x^m$,
definiujemy
\begin{equation}\label{eq:calka-lokalna}
 \int_M\omega:=\int_{x(U)} f(x^{-1}(u))\,\dd u^1\cdots\dd u^m.
\end{equation}
Dla dowolnego zwartego nośnika bierzemy podział
jedności $(\rho_j)$ z Twierdzenia
\ref{thm:podzial-jednosci-formy}, w mapach zgodnych
z orientacją, i kładziemy
\begin{equation}\label{eq:calka-globalna}
 \int_M\omega:=\sum_j\int_M\rho_j\omega.
\end{equation}
Tylko skończenie wiele składników może być niezerowych:
lokalnie skończona rodzina nośników spotyka zwarty
$\supp\omega$ tylko dla skończenie wielu indeksów.
W wymiarze $0$ definicję zastępuje suma
$\int_M f=\sum_{p\in M}\varepsilon_p f(p)$,
gdzie $\varepsilon_p\in\{1,-1\}$ jest orientacją punktu.
Zwarty nośnik w przestrzeni dyskretnej jest skończony.
\end{definicja}

\begin{twierdzenie}[Niezależność od map i podziału]
\label{thm:calka-dobrze-okreslona}
Definicja~\ref{def:calka-z-formy} nie zależy od
wybranego dodatniego atlasu ani podziału jedności.
Jeżeli $F:M\to N$ jest dyfeomorfizmem zachowującym
orientację, to dla $\omega\in\Omega_c^m(N)$ zachodzi
$\int_M F^*\omega=\int_N\omega$; gdy $F$ odwraca
orientację i $M,N$ są spójne, po prawej pojawia się
znak minus.
\end{twierdzenie}
\begin{proof}
W wymiarze $0$ teza wynika z sumy ze znakami i z tego,
że dyfeomorfizm jest bijekcją punktów. Dla $m\geq1$
najpierw niech nośnik leży w dwóch dodatnich mapach
$x$ i $y$. Z równości
\[
 \dd y^1\wedge\cdots\wedge\dd y^m
 =\det D(y\circ x^{-1})\,
   \dd x^1\wedge\cdots\wedge\dd x^m
\]
wynika, że współczynniki opisujące tę samą formę
w obu mapach różnią się o wyznacznik Jacobiego.
Jest on dodatni. Twierdzenie o zmianie zmiennych
w całce na otwartych podzbiorach $\R^m$ pokazuje,
że oba wzory~\eqref{eq:calka-lokalna} dają tę samą
liczbę. Dla mapy odwracającej orientację wyznacznik
jest ujemny, a zwykły wzór na zmianę zmiennych
zawiera jego wartość bezwzględną; stąd minus.

Jeśli $(\rho_j)$ i $(\sigma_k)$ są dwoma podziałami,
to na nośniku $\omega$ możemy wstawić
$1=\sum_k\sigma_k$ do każdej całki $\rho_j\omega$.
Otrzymujemy sumę $\sum_{j,k}\int_M\rho_j\sigma_k\omega$.
Każdy składnik ma nośnik w przecięciu dwóch map.
Rozcinamy go w razie potrzeby na mniejsze mapy,
stosujemy poprzedni rachunek i zamieniamy kolejność
obu skończonych sum. Wynik jest taki sam, gdy
zaczniemy od $(\sigma_k)$. Wzór dla $F$ sprawdzamy
na składnikach lokalnych przez $y\circ F\circ x^{-1}$
i sumujemy.
\end{proof}

Na zwartej rozmaitości każda gładka forma ma zwarty
nośnik. Dla zorientowanej krzywej parametryzowanej
$\gamma:[a,b]\to M$ całka z $1$-formy jest
$\int_\gamma\alpha:=\int_a^b\gamma^*\alpha$;
przeparametryzowanie zachowujące kierunek nie zmienia
jej wartości. Dla parametryzowanej powierzchni
analogicznie cofamy formę $2$-stopnia.

\begin{przyklad}[Całka po okręgu]
Dla $\gamma(t)=(\cos t,\sin t)$, $0\leq t\leq2\pi$,
zorientowanej przeciwnie do ruchu wskazówek zegara,
forma $\alpha=x\,\dd y-y\,\dd x$ daje
$\gamma^*\alpha=(\cos^2t+\sin^2t)\dd t=\dd t$.
Stąd $\int_{S^1}\alpha=2\pi$.
Nie jest to długość łuku wyprowadzona z metryki:
całkujemy $1$-formę wybraną w otaczającej płaszczyźnie.
\end{przyklad}

\section{Twierdzenie Stokesa}
\label{sec:formy-stokes}

Intuicja jest jedna: gdy zsumujemy brzegi
małych, zgodnie zorientowanych fragmentów,
ich wewnętrzne krawędzie wystąpią po dwa razy
z przeciwnymi znakami. Pozostanie brzeg całości.
Rachunek na formach pozwala zapisać tę zasadę
bez ograniczania się do konkretnych wymiarów.

\begin{figure}[!htbp]
\centering
\begin{tikzpicture}[>=Stealth,line cap=round,line join=round]
 \fill[manifoldblue!10] (0,0) rectangle (3,2.35);
 \fill[manifoldgreen!12] (3,0) rectangle (6,2.35);
 \draw[manifoldblue!70!black,thick] (0,0) rectangle (3,2.35);
 \draw[manifoldgreen!70!black,thick] (3,0) rectangle (6,2.35);
 \draw[->,manifoldred,very thick] (2.92,.6)--(2.92,1.68);
 \draw[->,manifoldgold!85!black,very thick] (3.08,1.68)--(3.08,.6);
 \draw[->,manifoldblue!70!black,thick] (.65,0)--(1.75,0);
 \draw[->,manifoldgreen!70!black,thick] (4.25,0)--(5.35,0);
 \node[manifoldblue!80!black] at (1.5,1.25) {$P_1$};
 \node[manifoldgreen!70!black] at (4.5,1.25) {$P_2$};
 \node[font=\small,align=left,anchor=west,text width=5.0cm]
  at (6.55,1.2)
  {Wspólna krawędź jest brzegiem $P_1$
   i $P_2$ z przeciwnymi kierunkami.
   W sumie znika; zostaje tylko brzeg zewnętrzny.};
\end{tikzpicture}
\caption{Znak orientacji usuwa wewnętrzną krawędź
przy składaniu dwóch fragmentów.}
\label{fig:stokes-anulowanie}
\end{figure}

\begin{twierdzenie}[Stokes dla form]
\label{thm:stokes-formy}
Niech $M^m$ będzie zorientowaną gładką rozmaitością
z brzegiem, $m\geq1$, a $\omega\in\Omega_c^{m-1}(M)$.
Brzeg ma orientację z Definicji
\ref{def:orientacja-brzegu}, a
$i:\partial M\hookrightarrow M$ jest \emph{inkluzją}.
Wtedy
\begin{equation}\label{eq:stokes-formy}
 \boxed{\displaystyle \int_M\dd\omega
        =\int_{\partial M}i^*\omega.}
\end{equation}
Gdy $\partial M=\varnothing$, prawa strona wynosi $0$.
\end{twierdzenie}

Zapis $i^*\omega$ oznacza \emph{cofnięcie formy}
przez inkluzję $i$, a nie nową pochodną czy
ograniczenie samego współczynnika. Z definicji
cofnięcia~\eqref{eq:pullback-formy}, dla
$p\in\partial M$ i $v_j\in T_p\partial M$,
\begin{equation}\label{eq:inkluzja-formy}
 (i^*\omega)_p(v_1,\ldots,v_{m-1})
 =\omega_p(\dd i_pv_1,\ldots,\dd i_pv_{m-1}).
\end{equation}
Różniczka $\dd i_p$ po prostu traktuje wektor
styczny do brzegu jako wektor styczny do $M$.
W mapie przybrzegowej $x^m\geq0$ mamy
$i(x^1,\ldots,x^{m-1})=(x^1,\ldots,x^{m-1},0)$,
więc $i^*(\dd x^j)=\dd x^j$ dla $j<m$, natomiast
$i^*(\dd x^m)=0$. Współczynniki formy podstawiamy
przy $x^m=0$; składniki zawierające $\dd x^m$
znikają. Znak orientacji brzegu rozpatrujemy
\emph{osobno}: nie jest częścią operacji $i^*$.

\begin{proof}
\emph{Przypadek mapy wewnętrznej.}
Załóżmy najpierw, że nośnik $\omega$ jest zwarty
i zawarty w dziedzinie mapy, która nie spotyka
$\partial M$. We współrzędnych jej obraz jest
otwartym podzbiorem $\R^m$. Ponieważ nośnik
leży wewnątrz obrazu, przedłużenie współczynników
przez zero na całą $\R^m$ jest gładkie i ma
zwarty nośnik. Dla standardowej orientacji piszemy
\[
 \omega=\sum_{i=1}^m f_i\,
 \dd x^1\wedge\cdots\wedge\widehat{\dd x^i}
      \wedge\cdots\wedge\dd x^m.
\]
Symbol daszka oznacza opuszczenie jednego czynnika.
Ze wzoru na różniczkę zewnętrzną wyraz
$\partial_j f_i\,\dd x^j$ mnożymy przez wszystkie
$\dd x^r$ z $r\ne i$. Jeśli $j\ne i$, występuje
dwukrotnie $\dd x^j$ i wyraz znika. Jeśli $j=i$,
przestawiamy $\dd x^i$ przez $i-1$ czynników,
co daje znak $(-1)^{i-1}$. Zatem
\[
 \dd\omega=\sum_{i=1}^m(-1)^{i-1}
 (\partial_i f_i)\,\dd x^1\wedge\cdots\wedge\dd x^m.
\]
Ustalmy wszystkie współrzędne poza $x^i$.
Funkcja $x^i\mapsto f_i(x)$ znika przy
$x^i\to\pm\infty$. Podstawowe twierdzenie
rachunku całkowego daje
\[
 \int_{-\infty}^{\infty}\partial_i f_i(x)\,\dd x^i
 =\lim_{b\to\infty}f_i(x^1,\ldots,b,\ldots,x^m)
  -\lim_{a\to-\infty}f_i(x^1,\ldots,a,\ldots,x^m)=0.
\]
Zamieniamy kolejność całkowania dzięki zwartemu
nośnikowi i sumujemy po $i$. Wychodzi
$\int_M\dd\omega=0$. Także
$i^*\omega=0$, ponieważ nośnik $\omega$ nie
spotyka brzegu, więc równość Stokesa zachodzi.

\emph{Przypadek mapy przybrzegowej.}
Teraz niech nośnik leży w mapie na otwarty
podzbiór półprzestrzeni
$\mathbb H^m=\{(x',s):x'\in\R^{m-1},\ s\geq0\}$,
gdzie $s=x^m$. Na razie używamy standardowej
orientacji $\dd x^1\wedge\cdots\wedge\dd x^m$.
Przedłużamy $f_i$ przez zero w obrębie
$\mathbb H^m$ poza obraz mapy; zwarty nośnik
sprawia, że to przedłużenie jest gładkie.
Nie przedłużamy funkcji przez $s=0$ na $s<0$:
wartość na brzegu może przecież być niezerowa.
Wzór na $\dd\omega$ jest taki sam jak wyżej.
Jeśli $i<m$, to po ustaleniu $s$ oraz pozostałych
współrzędnych stycznych całka z $\partial_i f_i$
po całej prostej $x^i$ wynosi zero.
Pozostaje wyraz $i=m$. Dla każdego $x'$ funkcja
$s\mapsto f_m(x',s)$ znika dla dostatecznie
dużego $s$. Stąd
\[
 \int_0^\infty\partial_s f_m(x',s)\,\dd s
 =\lim_{s\to\infty}f_m(x',s)-f_m(x',0)
 =-f_m(x',0).
\]
Po uwzględnieniu znaku $(-1)^{m-1}$ we wzorze
na $\dd\omega$ otrzymujemy
\[
 \int_{\mathbb H^m}\dd\omega
 =(-1)^{m-1}\int_{\R^{m-1}}
  \left(\int_0^\infty\partial_s f_m(x',s)\,\dd s\right)
          \dd x'
 =(-1)^m\int_{\R^{m-1}} f_m(x',0)\,\dd x'.
\]
Dlaczego jest to całka po \emph{zorientowanym}
brzegu? Na $s=0$ normalna na zewnątrz to
$n=-\partial_s$. Dla standardowej formy
objętościowej $\mu=\dd x^1\wedge\cdots\wedge\dd x^m$
liczymy bez skrótu
\[
 \mu(-\partial_s,\partial_1,\ldots,\partial_{m-1})
 =-(-1)^{m-1}=(-1)^m.
\]
Zasada „normalna na zewnątrz jako pierwsza”
z Definicji~\ref{def:orientacja-brzegu} mówi więc,
że orientacja brzegu względem standardowego
$\dd x^1\wedge\cdots\wedge\dd x^{m-1}$
ma znak $(-1)^m$. Jednocześnie wzór
\eqref{eq:inkluzja-formy} daje
\[
 i^*\omega=f_m(x',0)\,
 \dd x^1\wedge\cdots\wedge\dd x^{m-1},
\]
bo tylko wyraz z opuszczonym $\dd x^m$
nie zawiera różniczki normalnej. Całka tej
formy z orientacją brzegu jest zatem równa
obliczonej całce z $\dd\omega$.
Dla $m=1$ rozumiemy $x'$ jako pusty układ
współrzędnych: brzeg w tym modelu składa się
z punktu $s=0$ o znaku $-1$, a wzór nadal działa.
Gdy wybrana mapa odwraca orientację $M$,
całka po mapie i indukowana orientacja jej brzegu
zmieniają znaki jednocześnie, więc równość
pozostaje prawdziwa.

\emph{Przejście do całej rozmaitości.}
Wybierzmy podział jedności $(\rho_j)$
podporządkowany mapom wewnętrznym i przybrzegowym,
tak że każdy $\rho_j$ ma zwarty nośnik w swojej
mapie (Twierdzenie~\ref{thm:podzial-jednosci-formy}).
Ponieważ rodzina jest lokalnie skończona, a
$\supp\omega$ zwarty, tylko skończenie wiele
iloczynów $\rho_j\omega$ jest niezerowych.
Każdy ma nośnik zwarty w odpowiedniej mapie,
więc należy do jednego z dwóch powyższych
przypadków. Mamy $\sum_j\rho_j\omega=\omega$;
stosując liniowość $\dd$ do tej \emph{skończonej}
sumy, dostajemy $\sum_j\dd(\rho_j\omega)=\dd\omega$.
Z drugiej strony cofnięcie przez inkluzję jest
liniowe, a $i^*(\rho_j\omega)
=(\rho_j|_{\partial M})\,i^*\omega$.
Sumując lokalne równości Stokesa, otrzymujemy
\[
 \int_M\dd\omega
 =\sum_j\int_M\dd(\rho_j\omega)
 =\sum_j\int_{\partial M}i^*(\rho_j\omega)
 =\int_{\partial M}i^*\omega,
\]
ponieważ $\sum_j\rho_j|_{\partial M}=1$.
\end{proof}

\begin{przyklad}[Cofnięcie $1$-formy na brzeg dysku]
Niech $M=D^2$, $\omega=P(x,y)\,\dd x+Q(x,y)\,\dd y$
i niech $\gamma(t)=(\cos t,\sin t)$ parametryzuje
dodatnio zorientowany brzeg. Ponieważ
$\dd x(\gamma'(t))=-\sin t$ oraz
$\dd y(\gamma'(t))=\cos t$, dostajemy
\[
 \gamma^*(i^*\omega)
 =(i\circ\gamma)^*\omega
 =\bigl[-P(\cos t,\sin t)\sin t
       +Q(\cos t,\sin t)\cos t\bigr]\,\dd t.
\]
Tak obliczamy prawą stronę twierdzenia Stokesa
jako zwykłą całkę po $t\in[0,2\pi]$.
\end{przyklad}

\begin{wniosek}[Podstawowe twierdzenie rachunku całkowego]
\label{cor:stokes-ftc}
Jeśli $f$ jest gładka na $[a,b]$, to
\[
 \int_a^b f'(t)\,\dd t=f(b)-f(a).
\]
\end{wniosek}
\begin{proof}
Na zorientowanym przedziale $[a,b]$ forma $\dd f$
ma postać $f'(t)\dd t$. Jego brzeg ma orientację
$\partial[a,b]=\{b\}-\{a\}$, co sprawdziliśmy
w Stwierdzeniu~\ref{prop:brzeg-orientowalny}.
Całka $0$-formy $f$ po tym zorientowanym
zbiorze punktów jest z definicji $f(b)-f(a)$.
Wzór~\eqref{eq:stokes-formy} daje tezę.
\end{proof}
To identyfikacja jednowymiarowego przypadku Stokesa.
Podstawowego twierdzenia rachunku całkowego użyliśmy
już w jego dowodzie lokalnym, więc powyższy rachunek
nie stanowi niezależnego dowodu tego faktu analizy.

\begin{wniosek}[Twierdzenie Greena]
\label{cor:stokes-green}
Niech $D\subset\R^2$ będzie zwartym obszarem
o gładkim brzegu, z orientacją $\dd x\wedge\dd y$,
a $P,Q$ niech będą gładkie w otoczeniu $D$. Wtedy
\[
 \oint_{\partial D}P\,\dd x+Q\,\dd y
 =\iint_D\left(\frac{\partial Q}{\partial x}
             -\frac{\partial P}{\partial y}\right)
       \dd x\,\dd y.
\]
Orientacja $\partial D$ jest dodatnia, czyli
podczas obiegania każdej składowej brzegu
wnętrze $D$ pozostaje po lewej stronie.
\end{wniosek}
\begin{proof}
Dla $\alpha=P\,\dd x+Q\,\dd y$ reguła
Leibniza i $\dd x\wedge\dd x
=\dd y\wedge\dd y=0$ dają
\[
 \dd\alpha
 =P_y\,\dd y\wedge\dd x
  +Q_x\,\dd x\wedge\dd y
 =(Q_x-P_y)\,\dd x\wedge\dd y.
\]
Lewa strona wzoru Greena jest całką z
$i^*\alpha$ po brzegu; cofnięcie ogranicza
$\alpha$ do kierunków stycznych do krzywej,
jak w powyższym przykładzie z dyskiem. Podstawiamy $\alpha$ do
Twierdzenia~\ref{thm:stokes-formy}.
Znak orientacji sprawdzamy w punkcie
$(1,0)$ brzegu dodatnio zorientowanego dysku:
normalna zewnętrzna wskazuje w prawo, więc
dodatni wektor styczny wskazuje w górę.
Ta sama reguła działa na każdej składowej
brzegu, także na brzegu otworu.
\end{proof}

\begin{wniosek}[Klasyczne twierdzenie Stokesa o rotacji]
\label{cor:stokes-curl}
Niech $S\subset\R^3$ będzie zwartą gładką
zorientowaną powierzchnią z brzegiem,
$n$ jej dodatnią normalną jednostkową,
a $\mathbf F=(P,Q,R)$ gładkim polem w otoczeniu $S$.
Dla dodatniego kierunku $\partial S$ mamy
\[
 \oint_{\partial S}\mathbf F\cdot\dd\mathbf r
 =\iint_S(\nabla\times\mathbf F)\cdot n\,\dd A.
\]
\end{wniosek}
\begin{proof}
Bierzemy $1$-formę
$\alpha=P\,\dd x+Q\,\dd y+R\,\dd z$.
Bezpośredni rachunek pochodnych daje
\[
 \begin{aligned}
 \dd\alpha={}&(R_y-Q_z)\,\dd y\wedge\dd z\\
 &+(P_z-R_x)\,\dd z\wedge\dd x
 +(Q_x-P_y)\,\dd x\wedge\dd y.
 \end{aligned}
\]
Wektor współczynników przy trzech formach
bazowych to dokładnie
$\nabla\times\mathbf F
=(R_y-Q_z,P_z-R_x,Q_x-P_y)$.
Dla dodatniej parametryzacji
$r(u,v)$ fragmentu $S$ podstawmy
$a=r_u$, $b=r_v$. Z definicji iloczynu
zewnętrznego
\[
 (\dd\alpha)(a,b)
 =(\nabla\times\mathbf F)\cdot(a\times b)
 =(\nabla\times\mathbf F)\cdot n\,
    \lVert a\times b\rVert.
\]
Po scałkowaniu po $(u,v)$ otrzymujemy prawą
stronę tezy. Na krzywej brzegowej
$\alpha(\gamma')=\mathbf F(\gamma)\cdot\gamma'$,
więc całka z $i^*\alpha$ jest lewą stroną.
Twierdzenie~\ref{thm:stokes-formy} utożsamia
te dwie całki. Kierunek brzegu jest
indukowany przez orientację $S$; na płaskim
dysku z normalną w górę jest przeciwny
do ruchu wskazówek zegara.
\end{proof}

\begin{wniosek}[Twierdzenie Gaussa--Ostrogradskiego]
\label{cor:stokes-gauss}
Niech $D\subset\R^3$ będzie zwartym obszarem
o gładkim brzegu, a $\mathbf F=(P,Q,R)$
gładkim polem w otoczeniu $D$.
Jeśli $n$ oznacza normalną jednostkową
na zewnątrz, to
\[
 \iint_{\partial D}\mathbf F\cdot n\,\dd A
 =\iiint_D\operatorname{div}\mathbf F\,
                  \dd x\,\dd y\,\dd z.
\]
\end{wniosek}
\begin{proof}
Strumieniowi $\mathbf F$ odpowiada $2$-forma
\[
 \beta=P\,\dd y\wedge\dd z
      +Q\,\dd z\wedge\dd x
      +R\,\dd x\wedge\dd y.
\]
Różniczkując trzy składniki, dostajemy
\[
 \dd\beta=(P_x+Q_y+R_z)\,
           \dd x\wedge\dd y\wedge\dd z
          =(\operatorname{div}\mathbf F)\,
           \dd x\wedge\dd y\wedge\dd z.
\]
Dla dodatnio zorientowanej pary stycznych
wektorów $a,b\in T_p\partial D$ mamy
\[
 (i^*\beta)_p(a,b)
 =\mathbf F(p)\cdot(a\times b).
\]
Z reguły orientowania brzegu wynika, że
$a\times b$ wskazuje na zewnątrz; dla
ortonormalnej pary jest równy $n$.
Ogólnie jego długość dostarcza czynnika
elementu pola $\dd A$. Całka z $i^*\beta$
jest więc strumieniem przez brzeg, a całka
z $\dd\beta$ całką z dywergencji po $D$.
Równość daje Twierdzenie~\ref{thm:stokes-formy}.
\end{proof}

\begin{wniosek}[Zerowanie całki form zamkniętych i dokładnych]
\label{cor:stokes-top-form}
Jeśli $M^m$ jest zwartą zorientowaną rozmaitością
z brzegiem, a $\eta\in\Omega^{m-1}(M)$ jest
zamknięta, to $\int_{\partial M}i^*\eta=0$.
Jeżeli $M$ jest zwartą zorientowaną rozmaitością
bez brzegu i $\omega=\dd\eta\in\Omega^m(M)$,
to $\int_M\omega=0$. W szczególności forma
najwyższego stopnia o całce niezerowej nie jest
dokładna.
\end{wniosek}
\begin{proof}
W pierwszej sytuacji
$\int_{\partial M}i^*\eta=\int_M\dd\eta=0$.
W drugiej
$\int_M\omega=\int_{\partial M}i^*\eta=0$,
bo brzeg jest pusty. Oba kroki są
bezpośrednimi zastosowaniami~\eqref{eq:stokes-formy}.
\end{proof}
\section{Kohomologia de Rhama i ciąg Mayera--Vietorisa}
\label{sec:de-rham-mv}

\begin{definicja}[Kohomologia de Rhama]
\index{kohomologia de rhama@Kohomologia de Rhama}
\label{def:de-rham-cohomology}
\emph{Kompleks de Rhama} to ciąg
\[
 0\longrightarrow\Omega^0(M)\xrightarrow{\dd}
 \Omega^1(M)\xrightarrow{\dd}\cdots
 \xrightarrow{\dd}\Omega^m(M)\longrightarrow0.
\]
Jest kompleksem kołańcuchowym, bo $\dd^2=0$.
Formę z $\ker\dd$ nazywamy \emph{zamkniętą},
a formę z $\operatorname{im}\dd$ \emph{dokładną}. Definiujemy
\[
 H^k_{\mathrm{dR}}(M)
 :=\ker(\dd:\Omega^k\to\Omega^{k+1})/
   \operatorname{im}(\dd:\Omega^{k-1}\to\Omega^k).
\]
W stopniu $0$ jest to przestrzeń funkcji lokalnie stałych.
Z uwagi na brak form ujemnego stopnia zerowa funkcja
jest jedyną dokładną $0$-formą.
\end{definicja}

Iloczyn zewnętrzny przechodzi na klasy: jeśli
$\dd\alpha=\dd\beta=0$ i $\alpha$ zastąpimy przez
$\alpha+\dd\eta$, to różnica iloczynów wynosi
$(\dd\eta)\wedge\beta=\dd(\eta\wedge\beta)$.
Otrzymujemy \emph{pierścień kohomologii de Rhama}
$H^*_{\mathrm{dR}}(M)$ z przemiennością
$[\alpha][\beta]=(-1)^{k\ell}[\beta][\alpha]$.
Cofnięcie $F^*$ daje homomorfizm pierścieni
$H^*_{\mathrm{dR}}(N)\to H^*_{\mathrm{dR}}(M)$;
Wniosek~\ref{cor:de-rham-homotopia} pokazuje, że
gładkie homotopie dają tę samą operację.

\begin{twierdzenie}[Mayer--Vietoris dla form]
\label{thm:mv-de-rham}
Jeśli $M=U\cup V$ dla zbiorów otwartych, to
w każdym stopniu krótki ciąg
\begin{equation}\label{eq:mv-formy-kompleks}
\begin{gathered}
0\to\Omega^k(M)\xrightarrow{\ r\ }
\Omega^k(U)\oplus\Omega^k(V)
\xrightarrow{\ s\ }\Omega^k(U\cap V)\to0,\\
r(\omega)=(\omega|_U,\omega|_V),\qquad
s(\alpha,\beta)=\beta|_{U\cap V}-\alpha|_{U\cap V}.
\end{gathered}
\end{equation}
jest dokładny i zgodny z $\dd$. W konsekwencji
istnieje długi ciąg dokładny
\[
 \cdots\to H^{k-1}_{\mathrm{dR}}(U\cap V)
 \xrightarrow{\delta}H^k_{\mathrm{dR}}(M)
 \to H^k_{\mathrm{dR}}(U)\oplus H^k_{\mathrm{dR}}(V)
 \to H^k_{\mathrm{dR}}(U\cap V)
 \xrightarrow{\delta}H^{k+1}_{\mathrm{dR}}(M)\to\cdots.
\]
\end{twierdzenie}
\begin{proof}
Jeżeli $r(\omega)=0$, forma znika na $U$ i $V$,
więc znika na $M$: $r$ jest iniektywne.
Para $(\alpha,\beta)$ należy do $\ker s$ dokładnie
wtedy, gdy jej składniki zgadzają się na przecięciu;
wówczas sklejają się w jedną gładką formę na $M$.
To dowodzi $\operatorname{im}r=\ker s$.

Aby dowieść surjektywności $s$, weźmy
$\gamma\in\Omega^k(U\cap V)$ i podział jedności
$\rho_U+\rho_V=1$ podporządkowany $(U,V)$.
Na $U$ kładziemy $\alpha=-\rho_V\gamma$ na
$U\cap V$ i przedłużamy przez $0$ na $U\setminus V$;
na $V$ kładziemy $\beta=\rho_U\gamma$
i podobnie przedłużamy przez $0$.
Przedłużenia są gładkie: nośnik $\rho_V$ nie
dochodzi lokalnie do punktów poza $V$, a nośnik
$\rho_U$ do punktów poza $U$. Na przecięciu
$s(\alpha,\beta)=(\rho_U+\rho_V)\gamma=\gamma$.
Ponieważ obcięcie komutuje z $\dd$, otrzymujemy
krótki ciąg dokładny \emph{kompleksów}.
Twierdzenie~\ref{thm:les-kohomologia} z
Rozdziału~\ref{ch:algebra-homologiczna} tworzy
z niego długi ciąg. Konkretnie, dla zamkniętej
$\gamma$ wybieramy powyższe $(\alpha,\beta)$;
formy $\dd\alpha$ i $\dd\beta$ zgadzają się
na przecięciu, bo $\dd\gamma=0$. Sklejona
forma $\eta$ na $M$ jest zamknięta i
$\delta[\gamma]=[\eta]$. Zmiana podniesienia
lub reprezentanta $[\gamma]$ zmienia $\eta$
o formę dokładną, dokładnie jak w dowodzie
długiego ciągu kohomologii.
\end{proof}

\begin{przyklad}[Okrąg i forma kąta]
\label{prz:de-rham-s1}
Na $\R^2\setminus\{0\}$ forma
\[
 \vartheta=\frac{-y\,\dd x+x\,\dd y}{x^2+y^2}
\]
jest zamknięta: po zastosowaniu wzoru na $\dd$
otrzymujemy $\dd\vartheta=0$. Na okręgu
$\gamma(t)=(\cos t,\sin t)$ mamy
$\gamma^*\vartheta=\dd t$ i
$\int_{S^1}\vartheta=2\pi$. Gdyby $\vartheta$
była dokładna, Stokes na zamkniętej krzywej
dałby $0$, sprzeczność. Lokalnie możemy
pisać $\vartheta=\dd\theta$, ale funkcja kąta
po pełnym obrocie zmienia się o $2\pi$ i nie
jest globalną funkcją rzeczywistą.

W istocie $H^1_{\mathrm{dR}}(S^1)\cong\R$,
z izomorfizmem $[\alpha]\mapsto\int_{S^1}\alpha$.
Dowód surjektywności daje $c\vartheta$.
Jeśli całka z $\alpha$ znika, zapiszmy
$\gamma^*\alpha=g(t)\dd t$ dla funkcji
$2\pi$-okresowej $g$ i połóżmy
$h(t)=\int_0^t g(s)\,\dd s$. Wtedy
$h(t+2\pi)-h(t)=\int_t^{t+2\pi}g(s)\,\dd s
=\int_0^{2\pi}g(s)\,\dd s=0$ dla każdego $t$.
Funkcja $h$ jest zatem gładka i okresowa, więc schodzi
do gładkiej funkcji na $S^1$, której różniczka jest $\alpha$.
To dowodzi iniektywności. Ponadto
$H^0_{\mathrm{dR}}(S^1)=\R$ i wyższe grupy znikają.
\end{przyklad}

\begin{figure}[!htbp]
\centering
\begin{tikzpicture}[>=Stealth,line cap=round,line join=round]
 \fill[manifoldblue!9] (0,0) circle(2.0);
 \draw[manifoldblue!65!black,thick] (0,0) circle(2.0);
 \draw[manifoldgold!85!black,line width=1.4pt] (0,0) circle(1.35);
 \foreach \ang in {5,50,95,140,185,230,275,320}{
  \draw[->,manifoldred,thick]
   (\ang:1.35) arc[start angle=\ang,end angle={\ang+24},radius=1.35];
 }
 \fill[white] (0,0) circle(2.4pt);
 \draw[manifoldred,thick] (-.09,-.09)--(.09,.09)
  (-.09,.09)--(.09,-.09);
 \node[font=\small,align=left,anchor=west,text width=5.7cm]
  at (2.65,.5)
  {$\dd\vartheta=0$ poza środkiem,
   ale $\int_{S^1}\vartheta=2\pi$.
   Funkcja kąta jest pierwotną tylko lokalnie:
   po obejściu dziury zmienia się o $2\pi$.};
\end{tikzpicture}
\caption{Zamknięta $1$-forma na płaszczyźnie bez początku
może nie mieć globalnej funkcji pierwotnej.
Krzyżykiem zaznaczono jedyny usunięty punkt; niebieski
dysk przedstawia fragment dziedziny.}
\label{fig:de-rham-kat}
\end{figure}

\begin{przyklad}[Sfery i torus]
\label{prz:de-rham-sfery-torus}
Dwie dziedziny map na $S^n$ otrzymane przez usunięcie
przeciwnych biegunów są ściągalne, a ich przecięcie
zwija się na $S^{n-1}$. Długi ciąg z
Twierdzenia~\ref{thm:mv-de-rham}, zaczynając od
obliczenia dla $S^1$, daje indukcyjnie
\[
 H^k_{\mathrm{dR}}(S^n)=
 \begin{cases}\R,&k=0,n,\\0,&\text{w pozostałych stopniach,}\end{cases}
 \qquad n\geq1.
\]
Dla $n\geq2$ dokładność w stopniu $0$ mówi,
że $H^0(S^n)=\R$ i $H^1(S^n)=0$:
przecięcie jest spójne, a mapa
$H^0(U)\oplus H^0(V)=\R^2\to H^0(U\cap V)=\R$
ma postać $(a,b)\mapsto b-a$, więc jest surjektywna;
w stopniach $k\geq2$ odwzorowanie łączące
przenosi $H^{k-1}(S^{n-1})$ na $H^k(S^n)$.
Jedyny niezerowy stopień dodatni to $k=n$.

Na $T^2=S^1\times S^1$ niech
$\alpha=\pr_1^*\vartheta$ i
$\beta=\pr_2^*\vartheta$.
Ich całki po dwóch podstawowych okręgach
wynoszą odpowiednio $(2\pi,0)$ i $(0,2\pi)$;
$\int_{T^2}\alpha\wedge\beta=(2\pi)^2$.
Klasy $[\alpha],[\beta]$ są liniowo niezależne w stopniu $1$,
a $[\alpha\wedge\beta]$ jest niezerowa w stopniu $2$.
W obu przypadkach wynika to ze Stokesa: forma dokładna
ma zerową całkę po każdym z tych cykli.
Twierdzenie de Rhama poniżej i kompleks komórkowy
torusa z Przykładu~\ref{prz:homologia-cw-torus}
pokazują, że dają całe
$H^1_{\mathrm{dR}}(T^2)\cong\R^2$ oraz
$H^2_{\mathrm{dR}}(T^2)\cong\R$.
\end{przyklad}

\section{Dlaczego formy obliczają kohomologię przestrzeni?}
\label{sec:twierdzenie-de-rhama}

Forma zamknięta może nie mieć pierwotnej właśnie wtedy,
gdy wykrywa globalny „otwór”. Aby porównać ten język
z homologią singularną z Rozdziału~\ref{ch:homologia-przestrzeni},
całkujemy formę po zorientowanych sympleksach.
Używamy tu współczynników rzeczywistych: całka przyjmuje
wartości w $\R$, a torsja całkowitoliczbowa znika po
przejściu do tych współczynników.

\begin{definicja}[Kohomologia singularna z wartościami w $\R$]
\label{def:singular-cohom-r}
Jeżeli $C_k(M;\R)$ jest przestrzenią skończonych
$\R$-liniowych kombinacji ciągłych sympleksów
$\sigma:\Delta^k\to M$, to
$C^k(M;\R)=\operatorname{Hom}_{\R}(C_k(M;\R),\R)$.
Kobrzeg $\delta:C^k\to C^{k+1}$ określa wzór
$(\delta c)(\sigma)=c(\partial\sigma)$.
Z $\partial^2=0$ wynika $\delta^2=0$; iloraz
$\ker\delta/\operatorname{im}\delta$ oznaczamy przez
$H^k_{\mathrm{sing}}(M;\R)$.
\end{definicja}

Sympleks $\sigma:\Delta^k\to M$ nazywamy \emph{gładkim},
gdy w mapach jego składowe mają gładkie przedłużenia
do otoczeń punktów $\Delta^k$ w jego otoczce afinicznej.
Przy brzegu $M$ przedłużamy składowe do $\R^m$;
nie wymagamy, aby przedłużenie poza sympleks przyjmowało
wartości w półprzestrzeni modelującej $M$.
Na przykład $t\mapsto t$ z $[0,1]$ do $[0,\infty)$
jest w tym sensie gładkie. Dla takiego $\sigma$ definiujemy
$I(\omega)(\sigma)=\int_{\Delta^k}\sigma^*\omega$.
Znak orientacji $\Delta^k$ jest określony przez
kolejność jego wierzchołków. Na sympleksie
$\Delta^{k+1}$ twierdzenie Stokesa daje
\begin{equation}\label{eq:integracja-lancuchowa}
 \begin{aligned}
 (\delta I\omega)(\sigma)
 &= I\omega(\partial\sigma)
  =\sum_{j=0}^{k+1}(-1)^j
      \int_{[v_0,\ldots,\widehat v_j,\ldots,v_{k+1}]}
          \sigma^*\omega \\
 &=\int_{\Delta^{k+1}}\dd(\sigma^*\omega)
  =\int_{\Delta^{k+1}}\sigma^*(\dd\omega)
  =I(\dd\omega)(\sigma).
 \end{aligned}
\end{equation}
Uzasadnijmy użycie Stokesa na sympleksie, który ma naroża.
W otoczeniu punktu należącego do $r$ ścian jego dziedzinę
opisują współrzędne
$[0,\infty)^r\times\R^{k+1-r}$: równania przecinających
się ścian są niezależne. Rozkładamy formę podziałem
jedności i stosujemy lokalny dowód Stokesa do każdego
składnika o zwartym nośniku. Dla składnika z pominiętym
$\dd x^i$ całkowanie po $x^i$ daje zero, gdy $x^i\in\R$,
a gdy $x^i\geq0$, daje ujemną wartość współczynnika
na $x^i=0$. Są to dokładnie wkłady ścian z orientacją
„normalna na zewnątrz najpierw”. Przecięcia dwóch ścian
mają miarę zero w każdej ścianie i nie tworzą dodatkowych
całek. Dla orientacji sympleksu zadanej kolejnością
wierzchołków orientacja ściany naprzeciw $v_j$ jest
$(-1)^j$ razy orientacja pozostałych wierzchołków:
przesunięcie pominiętego wierzchołka na pierwsze miejsce
wymaga $j$ zamian, a przypadek $j=0$ daje orientację
brzegową dodatnią. Otrzymujemy więc dokładnie znaki
w \eqref{eq:integracja-lancuchowa}, bez zastępowania
naroży brzegiem gładkim.
Równanie mówi, że „całka po brzegu” i „różniczka
formy” opisują ten sam rachunek.

\begin{twierdzenie}[de Rham]
\label{thm:de-rham}
Dla każdej gładkiej rozmaitości $M$ (również z brzegiem)
całkowanie po gładkich sympleksach indukuje naturalny
izomorfizm przestrzeni wektorowych
\begin{equation}\label{eq:de-rham-iso}
 H^k_{\mathrm{dR}}(M)\xrightarrow{\ \cong\ }
 H^k_{\mathrm{sing}}(M;\R),
 \qquad [\omega]\longmapsto
 \bigl[\sigma\longmapsto\textstyle\int_\sigma\omega\bigr].
\end{equation}
Zapis po prawej najpierw określa kołańcuch na
\emph{gładkich} sympleksach. Przenosimy jego klasę do
zwykłej kohomologii singularnej przez odwrotność
kanonicznego izomorfizmu ograniczenia kołańcuchów,
którego istnienie udowodnimy poniżej. Nie wybieramy
niezależnie wygładzenia każdego ciągłego sympleksu:
takie wybory nie musiałyby być zgodne z jego ścianami.
\end{twierdzenie}
\begin{proof}
Dowód rozdzielimy na rachunek lokalny i precyzyjną
zasadę sklejania kompleksów. Równocześnie uzasadnimy,
dlaczego wystarczają gładkie sympleksy. Nie będziemy
potrzebować ani triangulacji $M$, ani metryki.

\emph{1. Trzy kompleksy i prostokąty wypukłe.}
Niech $C_*^\infty(U;\R)$ oznacza podkompleks
generowany przez gładkie sympleksy, a
$C_\infty^*(U;\R)=\operatorname{Hom}_{\R}
(C_*^\infty(U;\R),\R)$ jego kompleks dualny.
Ograniczenie kołańcuchów i całkowanie dają naturalne
odwzorowania kompleksów
\[
 C^*(U;\R)\xrightarrow{\ r\ }C_\infty^*(U;\R)
 \xleftarrow{\ I\ }\Omega^*(U).
\]
Pierwsze komutuje z kobrzegiem, ponieważ ściany
gładkiego sympleksu są gładkie; drugie dzięki
\eqref{eq:integracja-lancuchowa}.

Gdy $U$ jest niepustym wypukłym zbiorem otwartym
w $\R^m$ lub względnie otwartym w
$\mathbb H^m=\{x^m\geq0\}$, wybierzmy $a\in U$.
Homotopia $H(t,x)=(1-t)a+tx$ pozostaje w $U$.
Lemat~\ref{lem:homotopia-form} oblicza kohomologię
form, a operator pryzmatowy z Twierdzenia
\ref{thm:homotopia-singularna} oblicza obie kohomologie
singularne. Pryzmat i jego triangulacja są określone
afinicznie w dziedzinie, więc zachowują gładkość
sympleksów w przyjętym sensie. Każdy z trzech
kompleksów ma zatem kohomologię $\R$ w stopniu $0$
i zero w stopniach dodatnich. W stopniu $0$ oba
odwzorowania są identycznością na funkcjach stałych.
W szczególności $r$ i $I$ indukują izomorfizmy.

\emph{2. Algebraiczna zasada sklejania.}
Dla otwartego pokrycia $(U_i)$ rozważmy trzy tabele
\[
 \begin{aligned}
 A^{p,q}&=\prod_{i_0<\cdots<i_p}\Omega^q(U_{i_0\cdots i_p}),\\
 B^{p,q}&=\prod_{i_0<\cdots<i_p}C_\infty^q(U_{i_0\cdots i_p};\R),\\
 E^{p,q}&=\prod_{i_0<\cdots<i_p}C^q(U_{i_0\cdots i_p};\R),
 \qquad U_{i_0\cdots i_p}=U_{i_0}\cap\cdots\cap U_{i_p}.
 \end{aligned}
\]
Puste przecięcia wnoszą zera. W poziomie działa
naprzemienna suma ograniczeń $\delta_{\mathcal U}$,
np. $(\delta_{\mathcal U}a)_{ij}=a_j|_{U_{ij}}-a_i|_{U_{ij}}$.
W pionie działa odpowiednio $\dd$ lub kobrzeg singularny
$d_v$. Kompleks całkowity ma w stopniu $n$ sumę
składników $p+q=n$ i różniczkę
$D=\delta_{\mathcal U}+(-1)^p d_v$ na kolumnie $p$.
Ograniczenia komutują z $d_v$, więc $D^2=0$.

Użyjemy dwóch elementarnych własności takich tabel.
Jeśli wszystkie kolumny są dokładne, kompleks całkowity
jest dokładny. Istotnie, dla kocyklu $z$ jego składnik
w kolumnie o najmniejszym indeksie $p$ jest pionowo
zamknięty. Z dokładności wybieramy jego pionową pierwotną,
uwzględniając znak $(-1)^p$, i odejmujemy jej $D$.
Usuwamy tę kolumnę, ewentualnie zmieniając następną.
Powtarzamy to, aż usuniemy wszystkie składniki.
Procedura kończy się: w danym stopniu całkowitym jest
skończenie wiele par $(p,q)$. W najniższym pionowym
stopniu dokładność oznacza zerowe jądro, jeśli brak
już stopnia na pierwotną. Ten sam argument, z zamianą
ról indeksów, działa dla dokładnych wierszy.

Wynika stąd również zasada porównania: odwzorowanie
tabel, które indukuje izomorfizm kohomologii w każdej
kolumnie, indukuje izomorfizm kohomologii całkowitej.
Aby to sprawdzić, tworzymy stożek odwzorowania $f$.
Dla kompleksów $P,Q$ ma on składniki
$Q^n\oplus P^{n+1}$ i różniczkę
$(b,a)\mapsto(d_Qb+f(a),-d_Pa)$.
Z krótkiego ciągu dokładnego
$0\to Q\to\operatorname{Cone}(f)\to P[1]\to0$
i jego ciągu kohomologii wynika, że stożek jest dokładny
wtedy i tylko wtedy, gdy $f$ indukuje izomorfizmy
kohomologii. Dla odwzorowania tabel $f:P\to Q$ bierzemy
w pozycji $(p,q)$ przestrzeń $Q^{p,q}\oplus P^{p,q+1}$,
z różniczką pionową
$(b,a)\mapsto(d_vb+(-1)^p f(a),-d_va)$
i poziomą $(b,a)\mapsto(\delta_{\mathcal U}b,-\delta_{\mathcal U}a)$.
Te różniczki komutują, a po zastosowaniu znaku $(-1)^p$
w totalizacji otrzymujemy dokładnie stożek odwzorowania
kompleksów całkowitych. Kolumna $p$ jest stożkiem
$(-1)^p f$, więc jest dokładna na mocy założenia.
Stosujemy teraz poprzedni argument usuwania kolumn.
Dopuszczony dodatkowy stopień pionowy $-1$ nie przeszkadza:
przekątne nadal mają skończoną długość. Analogicznie,
jeśli do dokładnych wierszy dopiszemy w kolumnie $-1$
kompleks globalny $P$, to jego odwzorowanie do kompleksu
całkowitego $Q$ jest izomorfizmem kohomologii.
Rozszerzona tabela w stopniu całkowitym $n$ ma składniki
$Q^n\oplus P^{n+1}$. W kolumnie $-1$ pionowa różniczka
ma znak minus, a pozioma jest augmentacją $f$;
różniczka całkowita to zatem
$(b,a)\mapsto(D_Qb+f(a),-d_Pa)$.
Jest to dokładnie stożek tego odwzorowania, a jego
dokładne wiersze pozwalają zastosować ten sam argument.

\emph{3. Dlaczego wiersze obliczają kompleksy globalne.}
W przypadku $A$ dopisujemy $\Omega^q(M)$ na początku
wiersza. Wybierzmy gładki podział jedności
$\sum_j\rho_j=1$ podporządkowany pokryciu. Dla rodziny
naprzemiennej określamy
\[
 (ha)_{i_0\cdots i_{p-1}}
   =\sum_j\rho_j a_{j i_0\cdots i_{p-1}}.
\]
Składniki przedłużamy przez zero; lokalna skończoność
nośników zapewnia gładkość sumy. Dla $p=0$ otrzymujemy
globalną formę $\sum_j\rho_j a_j$.
Rozwinięcie sumy daje
$h\delta_{\mathcal U}+\delta_{\mathcal U}h=\id$:
wyrazy z pominiętym $i_r$ znoszą się parami, a pozostałe
dają $(\sum_j\rho_j)a$. Na formach globalnych złożenie
ograniczenia i $h$ także jest identycznością.
Rozszerzone wiersze są więc dokładne.

Dla $B$ oraz $E$ dopisujemy na początku kołańcuchy
na \emph{małych} sympleksach, czyli na sympleksach
o obrazie zawartym w jednym z $U_i$. Dla konkretnego
sympleksu $\sigma$ zbiór
$J_\sigma=\{i:\sigma(\Delta^q)\subset U_i\}$ jest
niepusty. Poziomy ciąg jego wartości to rozszerzony
kompleks kołańcuchów pełnego sympleksu o wierzchołkach
$J_\sigma$. Wybranie $i(\sigma)\in J_\sigma$ i wstawienie
go na początek indeksów daje homotopię $h$ spełniającą
tę samą tożsamość. Wiersz jest iloczynem takich ciągów
po wszystkich małych sympleksach; pozostaje dokładny,
ponieważ pierwotne wybieramy współrzędna po współrzędnej.

Lemat~\ref{lem:male-lancuchy} mówi, że włączenie małych
łańcuchów indukuje izomorfizm homologii. Jego dowód
działa też dla gładkich łańcuchów: podziały i homotopie
podziału są złożeniami z odwzorowaniami afinicznymi.
Nad $\R$ ograniczenie dualnych kołańcuchów indukuje
zatem izomorfizm kohomologii. Uzasadnienie algebraiczne
nie wymaga skończonego wymiaru: wybieramy dopełnienia
brzegów w cyklach i cykli w łańcuchach; kompleks jest
sumą części o zerowej różniczce (homologii) i par,
na których różniczka jest izomorfizmem. Dualizacja
zachowuje ten rozkład i dokładność części z parami.
W konsekwencji tabele $A,B,E$ obliczają odpowiednio
$H^*_{\mathrm{dR}}(M)$, $H^*(C_\infty^*(M;\R))$
i $H^*_{\mathrm{sing}}(M;\R)$.
Wszystkie ograniczenia i augmentacje komutują z $I$ i $r$.

\emph{4. Dwa zastosowania sklejania.}
Najpierw niech $W$ będzie otwartym podzbiorem $\R^m$
lub względnie otwartym podzbiorem $\mathbb H^m$.
Pokrywamy go prostokątami otwartymi, odpowiednio
przeciętymi z $\mathbb H^m$, i zawartymi w $W$.
Możemy wybrać pokrycie lokalnie skończone: w konstrukcji
wyczerpania zwartymi zbiorami z dowodu Twierdzenia
\ref{thm:podzial-jednosci-formy} zastępujemy małe kule
małymi prostokątami. Każde niepuste skończone przecięcie
jest wypukłe. Krok 1 pokazuje, że $I:A\to B$ oraz
$r:E\to B$ indukują izomorfizmy w każdej kolumnie.
Iloczyny przestrzeni wektorowych zachowują tu jądra,
obrazy i dokładność, więc także iloczyn lokalnych
izomorfizmów jest izomorfizmem. Kroki 2 i 3 dowodzą
obu porównań dla \emph{każdego} takiego $W$, nie tylko
dla zbioru wypukłego.

Teraz pokryjmy dowolną rozmaitość $M$ lokalnie skończoną
rodziną dziedzin map. Każde ich skończone przecięcie
jest otwartym podzbiorem jednej z tych dziedzin,
a więc jest dyfeomorficzne z pewnym $W$ rozpatrzonym
w poprzednim akapicie. Nie musi być ściągalne ani spójne.
Porównania $I$ i $r$ są na nim już udowodnione.
Ponowne zastosowanie kroków 2 i 3 daje izomorfizmy
\[
 H^*_{\mathrm{sing}}(M;\R)
 \xrightarrow{\ H(r)\ }H^*(C_\infty^*(M;\R))
 \xleftarrow{\ H(I)\ }H^*_{\mathrm{dR}}(M).
\]
Szukany izomorfizm to $H(r)^{-1}H(I)$.
Oba odwzorowania komutują z przenoszeniem sympleksów
i cofnięciem form przez gładkie odwzorowania, więc
izomorfizm jest naturalny. Cały argument obejmuje
rozmaitości z brzegiem dzięki modelom w $\mathbb H^m$.

Warto też wyjaśnić interpretację przez cykle. Nad ciałem
dualność kompleksów opisana w kroku 3 identyfikuje
kohomologię z dualną przestrzenią homologii.
Z izomorfizmu $H(r)$ wynika zatem, że włączenie
$C_*^\infty(M;\R)\to C_*(M;\R)$ indukuje izomorfizm
homologii: odwzorowanie liniowe, którego dualne jest
izomorfizmem, samo jest izomorfizmem (niezerowe jądro
lub niezerowy iloraz wykryłby funkcjonał liniowy).
Każda klasa cyklu ma więc gładkiego reprezentanta,
a dwaj tacy reprezentanci tej samej klasy różnią się
brzegiem gładkiego łańcucha. Dla formy zamkniętej ich
całki są równe przez \eqref{eq:integracja-lancuchowa}.
To uzasadnia niezależność okresów od reprezentantów
bez dokonywania niezgodnych wyborów na ścianach.
\end{proof}

\begin{przyklad}[Co widzą okresy, a czego nie widzą]
Jeżeli $\omega$ jest zamkniętą $k$-formą, jej
\emph{okres} na cyklu $z$ to $\int_z\omega$.
Jest niezależny od wyboru reprezentanta klasy
homologii: jeżeli $z$ zmienimy o $\partial c$,
różnica całek to $\int_c\dd\omega=0$.
Zmiana $\omega$ o formę dokładną także nie zmienia
okresu, bo $\int_z\dd\eta=\int_{\partial z}\eta=0$.
Na okręgu forma $\vartheta/(2\pi)$ z
Przykładu~\ref{prz:de-rham-s1} ma okres $1$ na
podstawowym cyklu. W $\mathbb{RP}^2$ grupa
$H_1(-;\Z)\cong\Z/2$ ma torsję, lecz
$H^1_{\mathrm{dR}}(\mathbb{RP}^2)=0$:
homomorfizm z $\Z/2$ do $\R$ musi wysłać
generator na liczbę $a$ z $2a=0$, zatem $a=0$.
Twierdzenie de Rhama wykrywa kohomologię
rzeczywistą; nie zastępuje informacji o torsji.
\end{przyklad}

\section{Porównanie homologii, kohomologii de Rhama i homotopii}
\label{sec:porownanie-trzech-teorii}

Pierwsza różnica dotyczy \emph{przestrzeni, na których
można zbudować daną teorię}. Homologię symplicjalną
definiujemy dla kompleksu symplicjalnego $K$:
sklejamy wierzchołki, odcinki, trójkąty i sympleksy
wyższych wymiarów wzdłuż ich ścian. Taki kompleks
\emph{nie musi być rozmaitością}. Na przykład dwa
wypełnione trójkąty stykające się tylko w jednym
wierzchołku tworzą kompleks, ale otoczenie
wspólnego wierzchołka nie przypomina dysku
ani półdysku: po usunięciu tego wierzchołka ma dwie
składowe, podczas gdy dysk oraz półdysk po usunięciu
punktu środkowego, odpowiednio punktu na prostym brzegu,
pozostają spójne. Mimo to łańcuchy, operator brzegu
$\partial$ i grupy $H_k^{\mathrm{simp}}(K;\Z)$
są określone.

\begin{figure}[!htbp]
\centering
\begin{tikzpicture}[line cap=round,line join=round]
 \fill[manifoldblue!14] (-3.1,-.95)--(-3.1,.95)--(0,0)--cycle;
 \fill[manifoldgold!19] (0,0)--(3.1,.95)--(3.1,-.95)--cycle;
 \draw[manifoldblue!75!black,thick]
  (-3.1,-.95)--(-3.1,.95)--(0,0)--cycle;
 \draw[manifoldgold!80!black,thick]
  (0,0)--(3.1,.95)--(3.1,-.95)--cycle;
 \fill[manifoldred] (0,0) circle (2.5pt);
 \node[below=5pt,font=\small] at (0,0) {$p$};
 \node[font=\small,manifoldblue!80!black] at (-2.05,0) {$\Delta^2$};
 \node[font=\small,manifoldgold!80!black] at (2.05,0) {$\Delta^2$};
\end{tikzpicture}
\caption{Dwa trójkąty ze wspólnym wierzchołkiem $p$.
To kompleks symplicjalny, ale w $p$ nie jest
powierzchnią gładką, nawet z brzegiem.}
\label{fig:kompleks-nie-rozmaitosc}
\end{figure}

Homologia singularna idzie jeszcze dalej: jej
sympleksy są \emph{ciągłymi odwzorowaniami}
$\Delta^k\to X$, dlatego definiuje się ją
dla dowolnej przestrzeni topologicznej $X$,
również bez wybranej triangulacji. Kohomologia
de Rhama wymaga natomiast \emph{gładkiej
rozmaitości} $M$: żeby pisać formy i ich
różniczkę $\dd$, potrzebujemy przestrzeni
stycznych i gładkich odwzorowań. Dla kompleksu
z Rysunku~\ref{fig:kompleks-nie-rozmaitosc}
zwykła kohomologia de Rhama w sensie
Definicji~\ref{def:de-rham-cohomology}
nie jest zdefiniowana. Grupy homotopii
można z kolei definiować dla przestrzeni
topologicznych z wybranym punktem bazowym.

Do czego służą te niezmienniki, jeśli słowo
„dziura” nie daje jeszcze ścisłej odpowiedzi?
Pozwalają rozstrzygać, czy cykl jest brzegiem,
czy zamknięta forma jest różniczką innej formy,
czy pętlę albo odwzorowanie sfery można
ściągnąć do punktu. Pomagają też odróżnić
przestrzenie oraz badać odwzorowania między nimi.
Istotny jest zarówno obiekt, od którego
zaczynamy, jak i relacja uznająca dwa takie
obiekty za równoważne:

\begin{center}
\small
\renewcommand{\arraystretch}{1.18}
\begin{tabular}{>{\raggedright\arraybackslash}p{3.05cm}
                >{\raggedright\arraybackslash}p{4.4cm}
                >{\raggedright\arraybackslash}p{6.75cm}}
\hline
\textbf{Typ}&\textbf{Na czym konstruowana}&
\textbf{Co wykrywa i o czym mówi}\\
\hline
$H_k^{\mathrm{simp}}(K;\Z)$,
homologia symplicjalna&
Zorientowane sympleksy kompleksu $K$,
całkowite łańcuchy, operator brzegu
$\partial$; cykle modulo brzegi.&
Rozpoznaje składowe ($H_0$), klasy cykli,
które nie są brzegami, ich liczbę niezależnych
generatorów i torsję. Na zamkniętej spójnej
rozmaitości jej najwyższa homologia z $\Z$
pomaga rozpoznać orientowalność.\\
\hline
$H_k^{\mathrm{sing}}(X;\Z)$,
homologia singularna&
Wszystkie ciągłe odwzorowania
$\Delta^k\to X$ jako generatory łańcuchów;
ten sam operator $\partial$ i iloraz
cykli przez brzegi.&
Daje homologiczne informacje bez wybierania
triangulacji, dla dowolnej przestrzeni
topologicznej. Ciągłe $f:X\to Y$ wyznacza
$f_*:H_k(X)\to H_k(Y)$; można nim badać,
jak odwzorowanie przenosi klasy cykli.\\
\hline
$H^k_{\mathrm{dR}}(M)$,
kohomologia de Rhama&
Gładkie formy różniczkowe na rozmaitości $M$,
różniczka $\dd$; formy zamknięte modulo dokładne.&
Mówi, czy zamknięta forma ma \emph{globalnie}
formę pierwotną. Całki po cyklach mierzą jej
okresy; iloczyn zewnętrzny form pozwala badać
iloczyny klas. Zwykła wersja ma wartości
rzeczywiste, więc nie wykrywa torsji.\\
\hline
$\pi_1(X,p)$,
grupa podstawowa&
Pętle o początku i końcu w $p$;
utożsamiamy pętle homotopijne, a działaniem
jest ich składanie.&
Mówi, których pętli nie można ściągnąć do
punktu. Zawiera informację o nakryciach
i może być nieprzemienna; przez to odróżnia
przestrzenie o tym samym $H_1$.\\
\hline
$\pi_k(X,p)$, $k\geq2$,
wyższe grupy homotopii&
Odwzorowania sfer $S^k\to X$ z punktem
bazowym, modulo homotopia zachowująca
ten punkt.&
Mówią, czy odwzorowanie sfery można
ściągnąć do punktu (równoważnie: przedłużyć
na dysk $D^{k+1}$). Mogą być niezerowe,
nawet gdy $H_k(X;\Z)=0$.\\
\hline
\end{tabular}
\end{center}

Na przykład $H_1$ pozwala wykazać, że
nie istnieje retrakcja dysku $D^2$ na jego
brzeg $S^1$: gdyby $r:D^2\to S^1$
i inkluzja $i:S^1\to D^2$ spełniały
$r\circ i=\id$, to na homologii byłoby
$r_*\circ i_*=\id$ na $H_1(S^1;\Z)\cong\Z$.
Jest to niemożliwe, bo $H_1(D^2;\Z)=0$.
Kohomologia de Rhama rozstrzyga, czy da się
zbudować globalny potencjał dla zamkniętej
formy, a $\pi_1$ pomaga opisać nakrycia.
„Dziury” są więc tylko jednym z obrazowych
skrótów dla tych konkretnych pytań.

Różni się też koszt rachunku. Bez triangulacji
kompleks łańcuchów singularnych ma ogromnie
wiele generatorów. Skończona triangulacja
sprowadza homologię symplicjalną do macierzy
operatorów brzegu, lecz znalezienie odpowiedniej
triangulacji i obliczenie jąder oraz obrazów
bywa trudne. Formy pozwalają czasem wykorzystać
współrzędne i symetrie rozmaitości, choć
ustalenie, czy forma ma \emph{globalną}
formę pierwotną, także wymaga pracy.
Nie ma jednej metody zawsze najłatwiejszej.

\subsection*{Od sympleksów do całek}

W tej części $H_k^{\mathrm{simp}}$ oznacza tę samą homologię,
którą wcześniej zapisywaliśmy $H_k^\Delta$, a
$H_k^{\mathrm{sing}}$ to wcześniejsze $H_k$ przestrzeni.
Indeksy służą tylko odróżnieniu konstrukcji w porównaniu.
Dla triangulacji $K$ Twierdzenie
\ref{thm:symplicjalna-singularna} utożsamia
$H_k^{\mathrm{simp}}(K;\Z)$ z
$H_k^{\mathrm{sing}}(M;\Z)$: wynik nie zależy
od wyboru triangulacji, chociaż liczba i kształt
sympleksów od niej zależą. Dla nieskończonej, lokalnie
skończonej triangulacji stosujemy ten sam wynik do skończonych
podkompleksów: skończony łańcuch symplicjalny ma skończony
nośnik, a obraz każdego łańcucha singularnego jest zwarty,
więc spotyka tylko skończenie wiele sympleksów.
To dotyczy także łańcuchów wypełniających; izomorfizmy
na skończonych podkompleksach dają więc zarówno
surjektywność, jak i iniektywność porównania globalnego.
Cyklem jest łańcuch
$z$ z $\partial z=0$. Jeżeli $z=\partial c$,
to reprezentuje zero w homologii.
W najprostszym przykładzie trzy krawędzie
obwodu trójkąta tworzą $1$-cykl. Jeśli
pozostawimy sam obwód, nie ma $2$-sympleksu,
którego ten cykl byłby brzegiem, i
$H_1\cong\Z$. Jeśli dodamy wypełniony trójkąt,
ten sam cykl staje się jego brzegiem:
\[
 \partial[v_0v_1v_2]
 =[v_1v_2]-[v_0v_2]+[v_0v_1].
\]
Wtedy $H_1=0$. To dosłowny sens
„wykrycia dziury” przez homologię.
Na okręgu
jedno pełne obejście daje generator $H_1(S^1;\Z)
\cong\Z$; na $S^2$ cała zorientowana powierzchnia
daje generator $H_2(S^2;\Z)\cong\Z$.

Klasa zamkniętej formy przypisuje liczby
klasom cykli: całkujemy formę po cyklu.
Wybieramy przy tym gładkiego reprezentanta całej klasy
cyklu, co uzasadniono w dowodzie Twierdzenia~\ref{thm:de-rham}.
Taką liniową regułę przypisującą wektorom
liczby nazywamy \emph{funkcjonałem liniowym}.
Właśnie w tym sensie kohomologia de Rhama
wiąże się z homologią. Jeśli
$\dd\omega=0$ i $\partial z=0$, definiujemy
\begin{equation}\label{eq:parowanie-de-rham-homologia}
 \langle[\omega],[z]\rangle:=\int_z\omega.
\end{equation}
Zmiana $\omega$ o $\dd\eta$ zmienia całkę o
$\int_{\partial z}\eta=0$; zmiana $z$ o
$\partial c$ zmienia ją o
$\int_c\dd\omega=0$. Jest więc dobrze określona
na klasach po obu stronach. To dokładnie
rachunek Stokesa~\eqref{eq:integracja-lancuchowa}.

Czy \emph{każdą} funkcję liniową przypisującą
klasom homologii liczby rzeczywiste można
otrzymać przez całkowanie zamkniętej formy?
Twierdzenie de Rhama daje odpowiedź twierdzącą.
Aby zobaczyć, skąd bierze się ta odpowiedź,
pokażmy najpierw krok algebraiczny.
Dla kompleksu
łańcuchowego przestrzeni wektorowych nad $\R$
kołańcuch $c:C_k\to\R$ jest kocyklem
($\delta c=0$) dokładnie wtedy, gdy znika
na $B_k=\partial C_{k+1}$. Jego ograniczenie
do $Z_k=\ker(\partial:C_k\to C_{k-1})$ zależy
więc tylko od klasy w $H_k=Z_k/B_k$.
Jeżeli ten funkcjonał na $H_k$ jest zerowy,
to $c$ znika na całym $Z_k$. Definiujemy
na $B_{k-1}=\partial C_k$ funkcjonał
$b(\partial a):=c(a)$. Jest dobrze określony:
z $\partial a=\partial a'$ wynika
$a-a'\in Z_k$, a zatem $c(a)=c(a')$.
Rozszerzamy $b$ liniowo na $C_{k-1}$.
Wtedy $(\delta b)(a)=b(\partial a)=c(a)$,
czyli $c$ był kobrzegiem. Odwrotnie, każdy
funkcjonał $\lambda:H_k\to\R$ definiuje
funkcjonał na $Z_k$, który można liniowo
przedłużyć na $C_k$; ponieważ znika na
$B_k$, przedłużenie jest kocyklem.
Dowiedliśmy zatem
\[
 H^k_{\mathrm{sing}}(M;\R)
 \cong \operatorname{Hom}_{\R}
       (H_k^{\mathrm{sing}}(M;\R),\R).
\]
Symbol $\operatorname{Hom}_{\R}(V,\R)$ oznacza
przestrzeń \emph{wszystkich} funkcjonałów liniowych
$V\to\R$: każdy przypisuje wektorowi z $V$
liczbę rzeczywistą i zachowuje dodawanie
oraz mnożenie przez skalar.
Łącząc to z Twierdzeniem~\ref{thm:de-rham}
i porównaniem symplicjalnym, otrzymujemy
dla triangulowanej rozmaitości
\begin{equation}\label{eq:porownanie-simp-de-rham}
 H^k_{\mathrm{dR}}(M)
 \cong \operatorname{Hom}_{\R}
       (H_k^{\mathrm{simp}}(K;\R),\R),
\end{equation}
przy czym izomorfizm jest dany przez całkę
z~\eqref{eq:parowanie-de-rham-homologia}.
Jeśli grupy homologii są skończenie generowane,
wymiar $H^k_{\mathrm{dR}}$ to $k$-ta liczba
Betti'ego, czyli liczba swobodnych generatorów
$H_k(M;\Z)$. Forma nie zobaczy torsji: gdy
$nz=\partial c$ dla $n\ne0$, to dla zamkniętej
$\omega$ mamy
$n\int_z\omega=\int_c\dd\omega=0$,
więc $\int_z\omega=0$.

Wybór współczynników wyjaśnia tu ważną różnicę.
Według Twierdzenia~\ref{thm:homologia-char-zero}
przejście z $\Z$ do $\Q$, $\R$ lub $\C$
zachowuje liczby Betti'ego, lecz usuwa torsję;
Przykład~\ref{prz:zmiana-wspolczynnikow-rp2-torus}
pokazuje to dla $\mathbb{RP}^2$ i torusa.
Zwykły kompleks form de Rhama jest rzeczywisty
i oblicza $H^k_{\mathrm{sing}}(M;\R)$.
Jeśli dopuścimy formy zespolone,
$\Omega^\bullet(M;\C)
=\Omega^\bullet(M;\R)\otimes_{\R}\C$,
otrzymamy analogicznie
$H^k_{\mathrm{dR}}(M;\C)
\cong H^k_{\mathrm{dR}}(M;\R)\otimes_{\R}\C$.
Współczynniki $\Q$ stosujemy bezpośrednio
w łańcuchach i kołańcuchach singularnych;
nie wymagają one osobnej teorii gładkich
form o wymiernych współczynnikach.

\subsection*{Od sfer do cykli}

W tym porównaniu zakładamy, że $M$ jest spójna drogowo;
w przeciwnym razie pracujemy na składowej zawierającej $p$.
Element $\pi_k(M,p)$ jest klasą odwzorowania
$(S^k,*)\to(M,p)$. Po przeniesieniu generatora
$H_k(S^k;\Z)$ otrzymujemy klasę w $H_k(M;\Z)$:
to odwzorowanie Hurewicza $h_k$
z Rozdziału~\ref{ch:wyzsza-homotopia}.
W stopniu $1$ dokładna odpowiedź z
Twierdzenia~\ref{thm:hurewicz-h1} brzmi
\[
 H_1(M;\Z)\cong\pi_1(M,p)_{\mathrm{ab}}.
\]
Komutatory pętli znikają po przejściu
do homologii, ponieważ dodawanie cykli
jest przemienne; $\pi_1$ może być nieprzemienna.
W wyższych stopniach $h_k$ nie musi być
izomorfizmem. Dla $k\geq2$, gdy $M$ jest
$(k-1)$-spójna, Twierdzenie
\ref{thm:hurewicz-wyzszy} utożsamia
pierwszą możliwą niezerową grupę
$\pi_k(M)$ z $H_k(M;\Z)$.
Aby zastosować tu wersję twierdzenia dla kompleksów CW,
używamy dodatkowego twierdzenia o triangulacji gładkich
rozmaitości: każda taka rozmaitość dopuszcza lokalnie
skończoną triangulację, a przy brzegu można wybrać
$\partial M$ jako podkompleks. Sympleksy tej triangulacji
są komórkami CW. Istnienie triangulacji jest osobnym
wynikiem topologii różniczkowej, nie konsekwencją de Rhama;
jego dowód wykracza poza to porównanie. Źródło wskazano
w dalszej lekturze. W rozważanych przykładach okręgu,
sfer, torusa i przestrzeni rzutowej strukturę CW znamy
już jawnie z wcześniejszych rozdziałów.
Teza dotyczy właśnie pierwszego
niezerowego stopnia przy podanej spójności;
nie utożsamia wszystkich późniejszych grup.

\begin{przyklad}[Ta sama informacja w stopniu pierwszym]
Na $S^1$ zachodzi
$\pi_1(S^1)\cong H_1(S^1;\Z)\cong\Z$,
a $H^1_{\mathrm{dR}}(S^1)\cong\R$.
Klasa formy $\vartheta/(2\pi)$ z
Przykładu~\ref{prz:de-rham-s1} daje wartość
$1$ na dodatnim generatorze homologii.
Liczba całkowita zlicza okrążenia, a całka
z formy przypisuje im liczby rzeczywiste
w sposób liniowy.
\end{przyklad}

\begin{przyklad}[Ta sama pierwsza homologia,
różne grupy podstawowe]
Torus $T^2$ ma $\pi_1(T^2)\cong\Z^2$.
Jeśli usuniemy otwarty dysk o domknięciu zawartym
we wnętrzu jego dwuwymiarowej komórki,
otrzymamy zwartą powierzchnię z brzegiem
$P=T^2\setminus\operatorname{int}D^2$.
W modelu komórkowym z Przykładu
\ref{prz:pi1-torus-rp2} część
dwuwymiarowa po usunięciu dysku ściąga się
na część jednowymiarową $S^1\vee S^1$.
Dlatego $\pi_1(P)\cong F_2$, wolna
nieprzemienna grupa na dwóch generatorach,
podczas gdy $H_1(P;\Z)\cong\Z^2$.
Zarówno $P$, jak i $T^2$ mają więc
$H^1_{\mathrm{dR}}\cong\R^2$:
w stopniu pierwszym formy i homologia
nie wykrywają nieprzemienności $\pi_1(P)$.
Pełne grupy homologii tych przestrzeni
różnią się jednak w stopniu $2$:
$H_2(P;\Z)=0$, a $H_2(T^2;\Z)=\Z$.
\end{przyklad}

\begin{przyklad}[Informacja sferyczna bez klasy homologii]
Dla $\mathbb{RP}^2$ wcześniejsze obliczenia dają
\[
 \pi_1(\mathbb{RP}^2)=\Z/2,\qquad
 \pi_2(\mathbb{RP}^2)=\Z,\qquad
 H_1(\mathbb{RP}^2;\Z)=\Z/2,\qquad
 H_2(\mathbb{RP}^2;\Z)=0.
\]
Porównaj Przykłady~\ref{prz:rp2-pi2}
i~\ref{prz:homologia-rp2}. Nakrycie
$S^2\to\mathbb{RP}^2$ pozostawia nietrywialną
klasę w $\pi_2$, ale jej obraz przez $h_2$
jest zerowy. Dodatkowo
$H^k_{\mathrm{dR}}(\mathbb{RP}^2)=0$
dla każdego $k>0$: całki form nie wykryją
ani torsji w $H_1$, ani tej klasy sferycznej.
Jest to wyraźny przykład, dlaczego trzech teorii
nie wolno po prostu utożsamiać.
\end{przyklad}

\section*{Dalsza lektura}
\addcontentsline{toc}{section}{Dalsza lektura}
\href{https://ocw.mit.edu/courses/18-965-geometry-of-manifolds-fall-2004/pages/lecture-notes/}
{Notatki kursu \emph{Geometry of Manifolds} (MIT OpenCourseWare)},
wykłady o formach, twierdzeniu Stokesa i twierdzeniu de Rhama.
Zob. także \href{https://www.math.toronto.edu/mgualt/courses/17-1300/docs/17-1300-notes-13.pdf}
{notatki o formach różniczkowych i twierdzeniu Stokesa}
Uniwersytetu w Toronto.
Porównanie z homologią i homotopią rozwija
\href{https://pi.math.cornell.edu/~hatcher/AT/AT.pdf}
{Allen Hatcher, \emph{Algebraic Topology}},
rozdziały 2--4. Twierdzenie o triangulacji rozmaitości
gładkich omawia Jacob Lurie w notatkach
\href{https://www.math.ias.edu/~lurie/937notes/937Lecture3.pdf}
{\emph{Whitehead Triangulations}, wykład 3} i kolejnych
wykładach tego kursu.
